\documentclass[11pt]{article}

% ----------------------------------------------------------------------
% Packages
% -----------------------------------------------------------------------
\usepackage[margin=1in]{geometry}
\usepackage{amsmath,amssymb,amsfonts,amsthm}
\usepackage{booktabs,array,tabularx,longtable,multirow}
\usepackage{subcaption,graphicx,float,placeins}
\usepackage{enumitem}
\usepackage{siunitx}
\usepackage{caption}
\usepackage{xcolor}
\definecolor{linkblue}{HTML}{1A73E8}
\usepackage{xspace}
\usepackage{pifont}
\usepackage{microtype}
\usepackage{natbib}
\usepackage[colorlinks=true, linkcolor=linkblue, citecolor=linkblue, urlcolor=linkblue]{hyperref}

% -----------------------------------------------------------------------
% Typesetting configuration
% -----------------------------------------------------------------------
\graphicspath{{figures/}}
\sisetup{detect-weight=true, detect-family=true}
\setlength{\tabcolsep}{4pt}
\renewcommand{\arraystretch}{1.12}
\captionsetup{font=small, labelfont=bf, width=0.95\textwidth,
              skip=4pt}
\setlength{\floatsep}{8pt plus 2pt minus 2pt}
\setlength{\textfloatsep}{10pt plus 2pt minus 4pt}
\setlength{\intextsep}{8pt plus 2pt minus 2pt}

% -----------------------------------------------------------------------
% Custom commands
% -----------------------------------------------------------------------
\newcommand{\RNAD}{RNAD\xspace}
\newcommand{\eg}{\textit{e.g.}\xspace}
\newcommand{\ie}{\textit{i.e.}\xspace}
\newcommand{\etal}{\textit{et~al.}\xspace}
\newcommand{\bm}[1]{\boldsymbol{#1}}
\DeclareMathOperator{\sgn}{sgn}
\newtheorem{remark}{Remark}

% Tighter paragraph spacing in lists
\setlist[itemize]{noitemsep, topsep=3pt, leftmargin=*}
\setlist[enumerate]{noitemsep, topsep=3pt, leftmargin=*, label=(\roman*)}

% -----------------------------------------------------------------------
\title{\RNAD: Regime-Normalized Adaptive Directional Loss\\
       for Financial Time-Series Forecasting}
\author{Nabeel Ahmad Saidd\\
Dr. A. P. J. Abdul Kalam Technical University\\
\texttt{nabeelahmadsaidd@gmail.com}}
\date{}

\begin{document}
\maketitle

% -----------------------------------------------------------------------
\begin{abstract}
Financial time-series forecasting faces a fundamental geometric
mismatch: standard training objectives minimize pointwise residuals on
an absolute price scale, yet financial returns exhibit time-varying
volatility, heavy tails, and regime switching---properties that
violate the homoskedastic, sub-Gaussian assumptions implicit in
mean-squared-error (MSE) training. This mismatch is acutely harmful
for one-step-ahead prediction, where both magnitude accuracy and
directional correctness jointly determine downstream decision quality.
We propose \textbf{\RNAD} (\textbf{R}egime-\textbf{N}ormalized
\textbf{A}daptive \textbf{D}irectional Loss), a drop-in training
objective that realigns optimization geometry with the structure of
financial prediction through three tightly integrated mechanisms:
(i)~\emph{volatility-adaptive normalization} via an exponential
moving average (EMA) blended with batch-wise scale estimates, ensuring
gradient magnitudes track predictive difficulty rather than price
amplitude; (ii)~\emph{tail-robust magnitude handling} via a
differentiable logistic gate that interpolates smoothly between a
quadratic core and Charbonnier-style sub-linear damping; and
(iii)~\emph{gated directional learning} that progressively downweights
sign supervision on economically small moves, suppressing gradient
noise in flat-price regimes. \RNAD is architecture-agnostic
and introduces minimal computational overhead: its sole statefulness
is a single per-asset scalar EMA volatility buffer.

We evaluate \RNAD comprehensively on nine financial assets spanning
equities, foreign exchange, cryptocurrencies, equity indices, and
commodities, across four modern forecasting architectures
(PatchTST, iTransformer, N-HiTS, TimeXer) under a strict
chronological evaluation protocol with no look-ahead bias. \RNAD
achieves the lowest RMSE on all four backbones and attains
the lowest per-asset RMSE on all nine PatchTST markets.
Controlled ablation studies indicate that each component contributes
independently: directional learning reduces RMSE by an average of
$0.70$ RMSE units across assets; tail gating improves robustness in
heavy-tailed cryptocurrency regimes; and noise gating exhibits modest,
regime-dependent effects. These results indicate
that principled loss-level design can yield forecasting improvements
of a comparable order of magnitude to several architectural choices,
without modifying model structure or incurring additional computational
cost at inference time.
\end{abstract}

% -----------------------------------------------------------------------
\section{Introduction}
\label{sec:intro}

Accurate financial time-series forecasting requires confronting a
structural misalignment between standard training objectives and the
statistical properties of market data. Asset returns exhibit
pronounced time-varying volatility~\citep{engle1982autoregressive,
bollerslev1986generalized}, heavy-tailed distributions that deviate
substantially from normality~\citep{mandelbrot1963variation,
cont2001empirical}, and intermittent regime
switching~\citep{hamilton1989new}. These characteristics violate the
homoskedastic, sub-Gaussian assumptions implicit in MSE training, which
treats every residual as if drawn from the same stationary distribution
regardless of the prevailing volatility environment.

A subtler but equally important deficiency of MSE is its
\emph{directional insensitivity}.
Many financial decisions---position entry, trend
following, mean-reversion triggers---are sign-sensitive: a model that
predicts the correct magnitude but the wrong direction is operationally
harmful, not merely inaccurate. Yet minimizing MSE imposes no
constraint on sign correctness, and a model can attain low MSE
while simultaneously producing directional signals that are no better
than random~\citep{leitch1991economic}. \citet{pesaran1992simple}
formally establish the statistical independence of directional accuracy
from conventional accuracy measures, providing direct motivation for
its incorporation into training objectives.

\paragraph{Source of the misalignment.}
The failure mode is structural and arises in two complementary
regimes. During high-volatility episodes, large shocks dominate
gradient updates and drive the optimizer to minimize absolute
magnitudes rather than to learn regime-relative behavior. During
low-volatility periods of negligible economic significance, directional
supervision injects irrelevant noise into the gradient signal.
This regime-blindness is particularly harmful in one-step-ahead
forecasting, where signal-to-noise ratios are intrinsically
low~\citep{cao2001financial,gu2020empirical}.

\paragraph{Loss-level design as a key lever.}
Rather than modifying model architectures---an approach that has been
extensively explored~\citep{vaswani2017attention,patchtst,
itransformer,challu2023nhits,zeng2023dlinear}---the optimization
objective itself can be redesigned to be regime-aware, tail-robust,
and sign-selective~\citep{barron2019general}. This approach preserves
full architectural flexibility while directly targeting the geometric
mismatch. We propose \RNAD (\textbf{R}egime-\textbf{N}ormalized
\textbf{A}daptive \textbf{D}irectional Loss), which unifies three
complementary mechanisms into a single differentiable training
objective:

\begin{enumerate}
\item \textbf{Adaptive Volatility Normalization.} Residuals and
  realized returns are normalized by an online volatility estimator
  (an EMA blended with the batch-wise root-mean-square return), so
  gradient magnitudes scale with predictive difficulty rather than
  with absolute price levels. This ensures that a single
  hyperparameter configuration behaves consistently across assets as
  disparate as BTCUSD (\$21{,}000) and AUDUSD (0.65).

\item \textbf{Tail-Robust Magnitude Term.} A differentiable logistic
  gate smoothly transitions the magnitude component from a quadratic
  kernel (preserving local gradient sensitivity for small,
  well-predicted errors) to Charbonnier-style sub-linear
  damping~\citep{charbonnier1994deterministic} (bounding gradient
  magnitudes during tail events), without the gradient discontinuity
  of hard thresholding as in Huber loss~\citep{huber1964robust}.

\item \textbf{Gated Directional Learning.} A smooth noise gate
  downweights sign-sensitive supervision when the realized move is
  economically small, suppressing directional-gradient noise in flat
  or near-flat regimes where the optimal sign boundary is ill-defined.
\end{enumerate}

\RNAD is a \emph{training-time intervention}, not a model modification.
Its sole statefulness is an $\mathcal{O}(1)$ per-asset scalar buffer,
and it is compatible with any gradient-based forecasting backbone.

\paragraph{Contributions.}
\begin{itemize}
\item We introduce \RNAD, a unified loss function combining adaptive
  volatility normalization, tail-robust magnitude handling, and gated
  directional learning. The formulation is mathematically explicit and
  operates entirely at the loss level, without modifying architectures,
  data pipelines, or inference procedures.

\item We conduct a large-scale benchmark spanning nine assets across
  five market classes, evaluating five loss functions on four modern
  architectures under a fixed chronological protocol, yielding
  20 backbone--loss combinations per asset evaluated under three
  seeds (180 unique configurations; 540 seed-level runs).

\item Controlled ablation studies isolate each component's contribution
  to both RMSE and directional accuracy, providing mechanistic
  insight beyond aggregate comparisons.

\item Diagnostic visualizations---volatility tracking, loss
  landscapes, gradient vector fields, and regime-stratified
  performance profiles---clarify how each loss component operates
  across different market conditions.
\end{itemize}

\paragraph{Paper organization.}
Section~\ref{sec:related} surveys related work and positions \RNAD
within the literature. Section~\ref{sec:method} formalizes the loss
function with precise notation. Section~\ref{sec:setup} specifies the
experimental protocol. Sections~\ref{sec:results}
and~\ref{sec:ablation} report empirical findings and ablation
evidence, respectively. Section~\ref{sec:discussion} interprets
the results. Section~\ref{sec:limitations} states the limitations
of the study. Section~\ref{sec:conclusion} concludes. Supplementary
diagnostics appear in the Appendix.

\FloatBarrier
% -----------------------------------------------------------------------
\section{Related Work}
\label{sec:related}

\subsection{Loss Functions for Regression and Time-Series Forecasting}

The squared-error loss (MSE) is the canonical regression criterion,
owing to its connection with the Gaussian log-likelihood and its
closed-form optimum, but its quadratic tail renders it sensitive to
outliers~\citep{huber1964robust}. The mean absolute error (MAE)
improves robustness at the cost of a non-differentiable origin and
uninformative zero gradients on well-fitted
samples~\citep{hastie2009elements}. Huber's combined loss interpolates
between L2 and L1 at a fixed transition threshold
$\delta$~\citep{huber1964robust}, trading sensitivity for stability;
log-cosh provides a smooth approximation to MAE with bounded second
derivatives~\citep{saleh2024logcosh}. Neither criterion adapts its
transition point to the prevailing data scale or volatility regime.

\citet{barron2019general} introduced a one-parameter family that
subsumes MSE, Cauchy, Welsch, and L1 losses under a single continuous
shape parameter, demonstrating that the loss shape can be tuned to
match the tail behavior of the noise distribution.
The Charbonnier penalty~\citep{charbonnier1994deterministic},
originally developed for edge-preserving image regularization, provides
the sub-linear asymptotic behavior employed here for outlier damping.
The tail-gating mechanism in \RNAD draws conceptual inspiration from
both works but operationalizes the quadratic-to-robust transition
via a \emph{runtime-adaptive} volatility signal rather than a fixed
hyperparameter.

In probabilistic forecasting, quantile regression
losses~\citep{koenker1978regression} and the asymmetric Laplace
likelihood enable heteroskedastic uncertainty estimation, but treat
magnitude and sign independently without directly incentivizing
directional correctness. Proper scoring rules, as surveyed by
\citet{gneiting2007strictly}, target calibration of distributional
forecasts rather than point-prediction accuracy. Recent work on
conformal prediction for time series~\citep{stankeviciute2021conformal}
extends calibration to arbitrary predictors, though also without
directional objectives. Pinball losses and quantile regression
forests~\citep{meinshausen2006quantile} remain standard in
risk-management applications but do not address regime-level
normalization.

\subsection{Directional Prediction in Financial Forecasting}

\citet{leitch1991economic} established that directional accuracy is an
economically motivated criterion that is statistically independent of
conventional error measures. \citet{pesaran1992simple} proposed a
nonparametric test for directional predictive ability, demonstrating
its independence from Diebold--Mariano-type forecast
comparisons~\citep{diebold1995comparing}. \citet{harvey1997testing}
extended these tests to handle non-normality and heteroskedasticity,
and \citet{clark2007approximately} proposed approximately normal
out-of-sample tests that remain valid under parameter estimation
uncertainty---a critical consideration in financial forecasting with
rolling windows. Together, these results motivate the direct
incorporation of sign-sensitive objectives into the training procedure.

Early work by \citet{white1988economic} and \citet{kuan1995forecasting}
pioneered neural-network applications to financial direction prediction,
relying on post-hoc accuracy evaluation. More recently,
\citet{michankow2024gmadl} proposed the Generalized Mean
Absolute Directional Loss (GMADL), which combines a magnitude
penalty with an explicit directional alignment term. However, GMADL
applies directional supervision unconditionally---regardless of the
economic significance of the realized move---and lacks volatility
normalization, making it susceptible to gradient noise during
low-volatility regimes and to gradient imbalance across assets with
disparate price levels. The gated directional component of \RNAD
addresses both deficiencies.

Directional objectives have also appeared outside finance. Focal
loss~\citep{focal} downweights easy negatives in object detection, a
structural analogue to the suppression of near-zero moves in \RNAD;
ordinal regression losses~\citep{gutierrez2016ordinal} encode rank
constraints that are related in spirit to sign constraints. In the
broader machine learning literature, pairwise ranking
losses~\citep{burges2005learning} and the maximum mean discrepancy
objective~\citep{gretton2012kernel} shape the output distribution in
ways analogous to directional gating, though without financial
regime-awareness.

\subsection{Regime-Aware and Adaptive Objectives}

Heteroskedasticity in financial returns has been extensively modeled
at the \emph{data} level via ARCH~\citep{engle1982autoregressive} and
GARCH~\citep{bollerslev1986generalized} families, which estimate
conditional variance to normalize returns. The volatility
normalization in \RNAD applies an analogous idea directly to the loss
geometry. RevIN~\citep{kim2022reversible} normalizes time-series inputs
via instance-wise statistics to mitigate non-stationarity and
distribution shift, establishing a precedent for normalization as a
preprocessing strategy within deep forecasters; \RNAD applies an
analogous normalization at the loss level, and the two approaches are
orthogonal and composable. Dish-TS~\citep{fan2023dish} and
SAN~\citep{liu2024adaptive} extend instance normalization to handle
within-sequence distribution shifts but do not address gradient-level
regime sensitivity.

Curriculum learning~\citep{bengio2009curriculum} and self-paced
learning~\citep{kumar2010self} dynamically adjust sample weights based
on training-phase difficulty, a concept related to the noise gating
in \RNAD, which suppresses supervision on ambiguous samples throughout
training. Multi-task learning with uncertainty
weighting~\citep{kendall2018multi} estimates per-task variances online
to balance gradient contributions; the per-asset EMA buffer in \RNAD
implements a simpler, single-task variant of this idea. Sample
weighting by prediction difficulty has also been studied in the
context of noisy-label learning~\citep{ren2018learning} and domain
adaptation~\citep{ganin2016domain}, where dynamically adjusting
gradient contributions from different sample subsets is central to
robust optimization.

\subsection{Neural Forecasting Architectures}

The encoder--decoder models based on dilated convolutions
(WaveNet;~\citealt{oord2016wavenet}) and autoregressive recurrent
networks (DeepAR;~\citealt{salinas2020deepar}) established canonical
sequence modeling baselines. N-BEATS~\citep{oreshkin2020nbeats} and
N-HiTS~\citep{challu2023nhits} subsequently demonstrated that
hierarchical decomposition stacks without recurrence could exceed
RNN-based methods on long-horizon benchmarks.

The transformer architecture of \citet{vaswani2017attention} has been
extensively adapted for time-series forecasting.
Informer~\citep{zhou2021informer} reduced attention complexity via
ProbSparse selection; Autoformer~\citep{wu2021autoformer} replaced
attention with auto-correlation decomposition; and
FEDformer~\citep{zhou2022fedformer} applied frequency-domain
transformations. PatchTST~\citep{patchtst} segments series into
patches and applies channel-independent self-attention with RevIN
normalization, achieving strong long-horizon performance.
iTransformer~\citep{itransformer} reverses the attention axis to
operate across variates rather than time steps, improving multivariate
generalization. Non-stationary Transformers~\citep{liu2022non} address
over-stationarization by reintegrating de-stationary statistics into
the attention mechanism, and Crossformer~\citep{zhang2023crossformer}
captures cross-dimensional dependencies via a two-stage routing
mechanism.

\citet{zeng2023dlinear} demonstrated that a single-layer linear model
can match transformer performance on several long-horizon benchmarks,
prompting the community to prioritize rigorous evaluation over
architectural complexity. TiDE~\citep{das2023long} revisited
MLP-based designs with separate encoder--decoder branches, recovering
competitive accuracy with reduced parameter counts.

More recent architectures include TimesNet~\citep{wu2023timesnet},
which projects time series into two-dimensional temporal--frequency
maps; TSMixer~\citep{ekambaram2023tsmixer}, applying MLP-mixer
operations; TimeMixer~\citep{wang2024timemixer}, combining multi-scale
decomposition and mixing; and ETSformer~\citep{woo2022etsformer},
which incorporates exponential smoothing into the attention mechanism.
Foundation models for time series---including
Chronos~\citep{ansari2024chronos}, TimesFM~\citep{das2024decoder},
and Time-LLM~\citep{jin2023time}---have further expanded the scope
by leveraging large-scale pretraining. These architectures serve as
evaluation contexts in the present study; no architectural novelty
is proposed here, and the experiments instead assess whether the
training-level intervention of \RNAD generalizes consistently across
this diverse family of backbones.

\subsection{Financial Forecasting Evaluation and Benchmarks}

Sound evaluation of financial forecasters requires chronological
splits, the absence of look-ahead bias, and replication across random
seeds to guard against initialization artifacts~\citep{makridakis2022m5,
zeng2023dlinear}. The Diebold--Mariano
test~\citep{diebold1995comparing} and the Model Confidence
Set~\citep{hansen2011model} provide statistical frameworks for
comparing forecast accuracy; \citet{diebold2015comparing} revisited
these in the context of modern machine learning models. Asset pricing
applications of machine learning~\citep{gu2020empirical,
freyberger2020dissecting,kelly2019characteristics} have demonstrated
that nonlinear models provide out-of-sample return predictability,
reinforcing the value of improving training objectives.
\citet{welch2008comprehensive} provide a comprehensive benchmark of
predictive variables for equity returns that highlights the difficulty
of out-of-sample improvement. \citet{campbell2008predicting} show
that economically motivated restrictions on forecasting models
improve their out-of-sample performance---an insight that the
directional gating of \RNAD embodies at the gradient level.

\paragraph{Positioning.}
Table~\ref{tab:loss-comparison} summarizes the property matrix of the
loss families discussed above. \RNAD is uniquely distinguished by
simultaneously satisfying all six desiderata: scale normalization,
tail robustness, directional awareness, online statefulness, noise
gating, and adaptive scaling. No prior loss function in the financial
forecasting literature satisfies all six criteria jointly.

\begin{table}[t]
\centering
\caption{Property matrix of loss functions. A checkmark indicates the
  property is satisfied by design; a cross indicates it is absent.
  \RNAD uniquely satisfies all six criteria. ``Fixed $\delta$''
  indicates that the transition is preset rather than data-adaptive.}
\label{tab:loss-comparison}
\small
\setlength{\tabcolsep}{3.5pt}
\begin{tabularx}{\textwidth}{X c c c c c c}
\toprule
Loss & Scale-Norm & Tail-Robust & Dir-Aware & Stateful
     & Noise-Gated & Adaptive-Scale \\
\midrule
MSE      & \ding{55} & \ding{55} & \ding{55} & \ding{55} & \ding{55} & \ding{55} \\
Huber    & \ding{55} & \ding{51}\ (fixed $\delta$) & \ding{55} & \ding{55} & \ding{55} & \ding{55} \\
log-cosh & \ding{55} & \ding{51}\ (smooth)         & \ding{55} & \ding{55} & \ding{55} & \ding{55} \\
GMADL    & \ding{55} & \ding{51}                   & \ding{51} & \ding{55} & \ding{55} & \ding{55} \\
\RNAD    & \ding{51} & \ding{51}\ (gated)          & \ding{51} & \ding{51} & \ding{51} & \ding{51} \\
\bottomrule
\end{tabularx}
\end{table}


Having established the position of \RNAD within the existing
literature, the following section develops the loss function
formally.

\FloatBarrier
% -----------------------------------------------------------------------
\section{The \RNAD Loss Function}
\label{sec:method}

\subsection{Notation and Problem Formulation}

Consider one-step-ahead forecasting of a univariate financial
close-price series $\{y_t\}_{t=1}^{T}$, where $y_t \in \mathbb{R}_{>0}$
denotes the close price at time $t$. A forecasting model
$f_\theta : \mathbb{R}^{d \times L} \to \mathbb{R}$ parameterized by
$\theta$ maps a window of $L$ past multivariate observations to a
predicted close price $\hat{y}_t = f_\theta(\mathbf{x}_{t-L:t})$.
In the implementation, both the input features and the target close
series are standardized with training-split $z$-score statistics
before the loss is computed; predictions are inverse-transformed only
for reporting raw-scale metrics.

For each time step $t$, three scalar quantities are defined:
\begin{align}
  e_t     &\;=\; \hat{y}_t - y_t,
          && \text{(forecast residual),}
  \label{eq:residual}\\
  r_t     &\;=\; y_t - y_{t-1},
          && \text{(realized one-step return),}
  \label{eq:realized}\\
  \hat{r}_t &\;=\; \hat{y}_t - y_{t-1},
          && \text{(predicted one-step return).}
  \label{eq:predicted-return}
\end{align}
Throughout, subscript $b \in \{1,\ldots,B\}$ indexes examples within
a mini-batch of size $B$. All three quantities in
Eqs.~\eqref{eq:residual}--\eqref{eq:predicted-return} are computable
from observations available at time $t$, ensuring the absence of
look-ahead bias.

\subsection{Adaptive Volatility Normalization}
\label{ssec:normalization}

A shared volatility scale $s_t > 0$ underlies all subsequent
computations. It blends two complementary estimators: a batch-wise
instantaneous estimate that responds promptly to volatility spikes,
and an EMA that retains regime-level memory across gradient steps.

\paragraph{Batch-wise scale.}
Given a mini-batch of realized returns $\{r_t^{(b)}\}_{b=1}^{B}$,
the batch-wise scale is defined as
\begin{equation}
  s_t^{\mathrm{batch}}
  = \sqrt{\frac{1}{B}\sum_{b=1}^{B}\bigl(r_t^{(b)}\bigr)^2 + \varepsilon},
  \qquad \varepsilon = 10^{-8}.
  \label{eq:batch-scale}
\end{equation}
This is the root-mean-square (RMS) of realized returns, which
approximates the conditional standard deviation under zero-mean
innovations. The regularizer $\varepsilon$ prevents division by zero
during sustained flat-price periods.

\paragraph{EMA scale.}
To capture regime-level volatility persistence, a running exponential
moving average is maintained:
\begin{equation}
  s_t^{\mathrm{EMA}}
  = \rho\,s_{t-1}^{\mathrm{EMA}}
  + (1-\rho)\,s_t^{\mathrm{batch}},
  \qquad \rho = 0.985,
  \label{eq:ema-scale}
\end{equation}
with initialization $s_0^{\mathrm{EMA}} = 1.0$. The effective memory
window spans approximately $1/(1-\rho) = 67$ gradient steps, which
is sufficient to track macro-level regime transitions while remaining
responsive to sustained volatility shifts.

\paragraph{Composite scale.}
The final scale is the geometric mean of the two estimates:
\begin{equation}
  s_t
  = \sqrt{s_t^{\mathrm{batch}} \cdot s_t^{\mathrm{EMA}} + \varepsilon}.
  \label{eq:composite-scale}
\end{equation}
Taking the geometric mean ensures that $s_t$ responds promptly to
volatility spikes via $s_t^{\mathrm{batch}}$ while being anchored to
the prevailing regime level via $s_t^{\mathrm{EMA}}$, without
tracking transient noise. The normalized quantities used throughout
are:
\begin{equation}
  u_t = \frac{e_t}{s_t},
  \qquad
  z_t = \frac{r_t}{s_t},
  \qquad
  \hat{z}_t = \frac{\hat{r}_t}{s_t},
  \label{eq:normalized}
\end{equation}
where $u_t$ is the volatility-normalized residual and $z_t$,
$\hat{z}_t$ are the normalized realized and predicted returns,
respectively.

\emph{Implementation note.} The codebase standardizes all input
features and the target close series using training-split statistics
before evaluating the loss; the notation above is written on that
standardized target scale, and raw prices are restored only for
reporting.

\begin{remark}
\normalfont
The EMA update in Eq.~\eqref{eq:ema-scale} is applied only during
training. During validation and testing, the loss is evaluated with
the EMA buffer frozen, so no evaluation-set statistics feed back into
the training state. Checkpoint-resume pipelines should therefore
persist the EMA buffer alongside model parameters.
\end{remark}

\subsection{Tail-Robust Magnitude Term}
\label{ssec:magnitude}

The magnitude component smoothly interpolates between a quadratic
kernel that preserves sensitivity for well-predicted samples and a
sub-linear Charbonnier kernel that bounds gradient magnitudes during
tail events.

\paragraph{Tail gate.}
Let $\sigma(x) = (1+e^{-x})^{-1}$ denote the logistic sigmoid. The
interpolation weight is
\begin{equation}
  g_t
  = \sigma\!\left(\frac{|u_t| - \tau_{\mathrm{mag}}}{\tau_{\mathrm{temp}}}\right),
  \label{eq:tail-gate}
\end{equation}
where $\tau_{\mathrm{mag}} = 1.1$ is the transition midpoint in units
of $s_t$ and $\tau_{\mathrm{temp}} = 0.3$ controls the sharpness of
the transition. When $|u_t| \ll \tau_{\mathrm{mag}}$, we have
$g_t \approx 0$ (quadratic regime); when $|u_t| \gg \tau_{\mathrm{mag}}$,
$g_t \approx 1$ (Charbonnier regime). The logistic transition avoids
the gradient discontinuity of the Huber threshold.

\paragraph{Kernel definitions.}
The quadratic and Charbonnier kernels are:
\begin{align}
  \ell_t^{\mathrm{L2}}
  &= \tfrac{1}{2}u_t^2,
  \label{eq:l2-kernel}\\[2pt]
  \ell_t^{\mathrm{charb}}
  &= \sqrt{u_t^2 + \varepsilon_c^2} - \varepsilon_c,
  \qquad \varepsilon_c = 10^{-3}.
  \label{eq:charb-kernel}
\end{align}
The offset $\varepsilon_c$ ensures differentiability at $u_t = 0$.
As $|u_t| \to \infty$, $\ell_t^{\mathrm{charb}} \approx |u_t|$
(linear growth), whereas $\ell_t^{\mathrm{L2}} = u_t^2$ grows
without bound, amplifying the influence of extreme residuals.

\paragraph{Blended magnitude loss.}
The final magnitude term is:
\begin{equation}
  \ell_t^{\mathrm{mag}}
  = (1 - g_t)\,\ell_t^{\mathrm{L2}}
  + g_t\,\ell_t^{\mathrm{charb}}.
  \label{eq:mag-loss}
\end{equation}
Because the gates depend on $|u_t|$ and $|z_t|$, the full objective
is piecewise smooth rather than globally $C^\infty$; it is
differentiable almost everywhere, with non-differentiable points only
at $u_t = 0$ and $z_t = 0$. Away from those measure-zero points,
the derivative with respect to $u_t$ transitions from approximately
$u_t$ near the origin to approximately $\mathrm{sgn}(u_t)$ in the
tails, bounding the influence of tail observations on parameter
updates.

\subsection{Directional Term with Noise Gating}
\label{ssec:directional}

The directional component penalizes sign disagreement between the
predicted and realized return, with a smooth gate that downweights
the penalty on economically small moves.

\paragraph{Smooth sign encoding.}
Both returns are mapped to $(-1, +1)$ via a scaled hyperbolic tangent:
\begin{equation}
  d_t^{\mathrm{true}} = \tanh(\beta\,z_t),
  \qquad
  d_t^{\mathrm{pred}} = \tanh(\beta\,\hat{z}_t),
  \qquad
  \beta = 4.5.
  \label{eq:tanh-sign}
\end{equation}
The sharpness parameter $\beta$ controls the transition width around
zero. At $\beta = 4.5$, a normalized move $|z_t| = 0.5$ is encoded
as $|\tanh(2.25)| \approx 0.978$, rendering the encoding near-binary
for moderate moves while remaining differentiable at the origin.

\paragraph{Directional alignment score.}
The per-sample alignment is:
\begin{equation}
  a_t
  = d_t^{\mathrm{true}} \cdot d_t^{\mathrm{pred}}
  \;\in\; [-1,\,1].
  \label{eq:dir-score}
\end{equation}
Perfect directional agreement ($\hat{z}_t$ and $z_t$ share sign)
gives $a_t \to +1$; perfect disagreement gives $a_t \to -1$.

\paragraph{Noise gate.}
A smooth sigmoid gate downweights directional supervision when the
realized move is negligible relative to the current volatility:
\begin{equation}
  w_t = \sigma\!\left(\frac{|z_t| - \tau_{\mathrm{noise}}}{\tau_{\mathrm{noise,temp}}}\right),
  \qquad
  \tau_{\mathrm{noise}} = 0.1.
  \label{eq:noise-gate}
\end{equation}
The gate temperature is fixed at $\tau_{\mathrm{noise,temp}} = 0.04$.
Both constants are expressed in normalized units and are held fixed
across all assets; no per-asset calibration is applied. Smaller moves
receive continuously smaller directional weights rather than a binary
on/off suppression.

\paragraph{Directional penalty.}
Using softplus for numerical stability, the directional loss is:
\begin{equation}
  \ell_t^{\mathrm{dir}}
  = w_t \cdot \mathrm{softplus}\!\left(-\kappa\,a_t\right),
  \qquad
  \mathrm{softplus}(x) = \ln(1 + e^x),
  \qquad
  \kappa = 4.5.
  \label{eq:dir-loss}
\end{equation}
For perfect alignment ($a_t = 1$): $\ell_t^{\mathrm{dir}} =
\mathrm{softplus}(-4.5) \approx 0.011$. For perfect disagreement
($a_t = -1$): $\ell_t^{\mathrm{dir}} = \mathrm{softplus}(4.5)
\approx 4.51$. This asymmetry strongly penalizes directional errors
while imposing negligible cost on correct predictions, creating a
well-defined decision boundary around the sign of $\hat{z}_t$.

\subsection{Full Objective}
\label{ssec:full-loss}

The per-sample \RNAD loss is:
\begin{equation}
  \mathcal{L}_t^{\RNAD}
  = \ell_t^{\mathrm{mag}}
  + \lambda_{\mathrm{dir}}\,\ell_t^{\mathrm{dir}},
  \qquad
  \lambda_{\mathrm{dir}} = 0.5,
  \label{eq:rnad-sample}
\end{equation}
and the mini-batch training objective averages over the batch:
\begin{equation}
  \mathcal{L}^{\RNAD}
  = \frac{1}{B}\sum_{b=1}^{B}\mathcal{L}_{b}^{\RNAD}.
  \label{eq:rnad-batch}
\end{equation}
The EMA buffer $s_t^{\mathrm{EMA}}$ is the sole persistent state;
it is updated after each gradient step. All other quantities in
Eqs.~\eqref{eq:tail-gate}--\eqref{eq:rnad-sample} are computed
from the current batch. Table~\ref{tab:hyperparams} lists all
loss-specific constants, and Section~\ref{sec:reproducibility}
summarizes the full backbone and training settings.
Fig.~\ref{fig:architecture} illustrates the full computational graph.

\FloatBarrier
\begin{figure}[!http]
\centering
\includegraphics[width=0.95\textwidth]{figure_01_rnad_schematic.pdf}
\caption{\textbf{\RNAD loss architecture (Simplified).}
  Model outputs $\hat{y}_t$ and target values $y_t$ enter two
  parallel pathways. \emph{Upper (magnitude) pathway:} the residual
  $e_t$ is normalized by the composite volatility scale $s_t$
  (Eq.~\ref{eq:composite-scale}) to yield $u_t$; the tail gate $g_t$
  (Eq.~\ref{eq:tail-gate}) interpolates between the quadratic kernel
  $\ell^{\mathrm{L2}}$ (Eq.~\ref{eq:l2-kernel}) and the Charbonnier
  kernel $\ell^{\mathrm{charb}}$ (Eq.~\ref{eq:charb-kernel}).
  \emph{Lower (directional) pathway:} normalized returns $z_t$ and
  $\hat{z}_t$ are sign-encoded via $\tanh$ (Eq.~\ref{eq:tanh-sign});
  the noise gate $w_t$ (Eq.~\ref{eq:noise-gate}) downweights
  supervision on small-magnitude moves; the directional penalty
  $\ell^{\mathrm{dir}}$ (Eq.~\ref{eq:dir-loss}) is computed via
  softplus. Both pathways are combined in
  Eq.~\eqref{eq:rnad-sample}. The EMA buffer $s_t^{\mathrm{EMA}}$
  (dashed arrow) is updated after each gradient step.}
\label{fig:architecture}
\end{figure}

\begin{figure}[!http]
\centering
\includegraphics[width=0.95\textwidth]{figure_02_rnad_surface.pdf}
\caption{\textbf{\RNAD loss surface.}
  $\mathcal{L}_t^{\RNAD}$ as a function of normalized residual $u_t$
  and normalized realized return $z_t$ (Eq.~\ref{eq:normalized}).
  Three geometric features are visible: (i)~a quadratic bowl for
  $|u_t| < \tau_{\mathrm{mag}} = 1.1$; (ii)~smooth Charbonnier
  flattening for $|u_t| > 1.1$; and (iii)~directional penalty ridges
  where $\mathrm{sgn}(\hat{z}_t) \neq \mathrm{sgn}(z_t)$ and
  $|z_t| > \tau_{\mathrm{noise}} = 0.1$. The surface is
  piecewise smooth and differentiable almost everywhere; the only
  non-differentiable points arise from the absolute-value gates at
  $u_t = 0$ and $z_t = 0$.}
\label{fig:rnad-surface}
\end{figure}

\FloatBarrier

\subsection{Mechanistic Rationale}
\label{ssec:rationale}

Each component of \RNAD targets a distinct failure mode of standard
losses on financial data:

\begin{itemize}
\item \textbf{Volatility normalization} (Eq.~\ref{eq:composite-scale})
  removes the spurious coupling between gradient magnitude and asset
  price level. Without normalization, high-priced assets
  (BTCUSD $\sim$\$21{,}000) generate gradients orders of magnitude
  larger than low-priced assets (AUDUSD $\sim$0.65), preventing a
  single hyperparameter configuration from being applicable across
  asset classes.

\item \textbf{Tail gating} (Eq.~\ref{eq:tail-gate}) prevents rare
  shock events from dominating optimization. Unlike Huber's hard
  threshold, the smooth transition preserves gradient informativeness
  within the interior of the residual distribution while bounding
  gradient norms in the tails, improving robustness without
  sacrificing sensitivity.

\item \textbf{Directional learning} (Eq.~\ref{eq:dir-loss})
  explicitly incentivizes sign correctness, addressing MSE's
  structural inability to enforce directional
  accuracy~\citep{leitch1991economic,pesaran1992simple}. As analyzed
  in Section~\ref{ssec:coupling}, this coupling between sign and
  magnitude objectives also reduces RMSE.

\item \textbf{Noise gating} (Eq.~\ref{eq:noise-gate}) focuses
  directional supervision on economically meaningful moves. Near-zero
  returns contribute noisy, potentially misleading sign gradients
  that can destabilize learning in low-volatility regimes such as
  foreign exchange markets (see ablation results in
  Table~\ref{tab:ablation}).
\end{itemize}

The experimental setup required to evaluate these mechanisms is
described in the following section.

\FloatBarrier
% -----------------------------------------------------------------------
\section{Experimental Setup}
\label{sec:setup}

\subsection{Datasets and Assets}

We evaluate on nine financial assets spanning five market classes
(Table~\ref{tab:data-summary}): equities (AAPL, DAX30, FTSE100),
foreign exchange (AUDUSD, USDJPY), cryptocurrencies (BTCUSD, ETHUSD),
equity indices (US500), and commodities (XAUUSD). All series consist
of minute-frequency OHLCV data. Total sample sizes range from
15{,}673 (AAPL) to 68{,}536 (USDJPY), spanning multiple market
regimes and volatility cycles. Realized volatility proxies
(training-set RMS returns) vary by two orders of magnitude across
assets, from USDJPY ($9.6 \times 10^{-4}$) to ETHUSD
($8.5 \times 10^{-3}$), providing strong motivation for the
cross-asset normalization mechanism.

\paragraph{Chronological split protocol.}
All splits are strictly chronological with no shuffling: $80\%$
training, $10\%$ validation, $10\%$ test. This protocol eliminates
look-ahead bias, preserves temporal autocorrelation structure, and
evaluates generalization to the most recent---and typically the most
distributional-shift-prone---period of each series, following
established best practice~\citep{makridakis2022m5,zeng2023dlinear,
welch2008comprehensive}.

\begin{table}[!htbp]
\centering
\caption{Dataset summary. All assets use an identical OHLCV and
  technical indicator feature set, a raw close-price target, and an
  $80/10/10$ chronological split. Vol.\ Proxy is the training-set
  root-mean-square one-step return, used to initialize
  $s_0^{\mathrm{EMA}}$.}
\label{tab:data-summary}
\small
\begin{tabularx}{\textwidth}{X l r r r r l}
\toprule
Asset   & Market Type    & Total     & Train      & Val       & Test      & Vol.\ Proxy \\
\midrule
AAPL    & Equity         & 15{,}673  & 12{,}538   & 1{,}567   & 1{,}568   & 0.00616 \\
AUDUSD  & Forex          & 68{,}535  & 54{,}828   & 6{,}853   & 6{,}854   & 0.00127 \\
BTCUSD  & Crypto         & 62{,}793  & 50{,}234   & 6{,}279   & 6{,}280   & 0.00653 \\
DAX30   & Equity Index   & 57{,}783  & 46{,}226   & 5{,}778   & 5{,}779   & 0.00233 \\
ETHUSD  & Crypto         & 62{,}783  & 50{,}226   & 6{,}278   & 6{,}279   & 0.00847 \\
FTSE100 & Equity Index   & 56{,}767  & 45{,}413   & 5{,}676   & 5{,}678   & 0.00197 \\
US500   & Equity Index   & 57{,}665  & 46{,}132   & 5{,}766   & 5{,}767   & 0.00200 \\
USDJPY  & Forex          & 68{,}536  & 54{,}828   & 6{,}853   & 6{,}855   & 0.00096 \\
XAUUSD  & Commodity      & 65{,}093  & 52{,}074   & 6{,}509   & 6{,}510   & 0.00164 \\
\bottomrule
\end{tabularx}
\end{table}

\FloatBarrier

\subsection{Feature Engineering and Preprocessing}

All models receive an identical, causally valid feature set computed
exclusively from observations available at or before time $t$:
raw OHLCV bars; MACD (line, signal, histogram); RSI-14; EMA-9; and
Bollinger Bands (SMA-20, upper band, lower band). The 13-column
feature vector is ordered as [open, high, low, close, volume, macd,
macd\_signal, macd\_hist, rsi\_14, ema\_9, bb\_middle, bb\_upper,
bb\_lower], so the forecast target is the close channel at index 3.
Both the input features and the close-price target are standardized
using $z$-score statistics ($\mu$, $\sigma$) computed on the training
split alone and applied identically to the validation and test splits.
The fitted target scaler is retained to inverse-transform predictions
before computing raw-scale RMSE, MAE, and directional accuracy. The
forecasting window is fixed at $L = 24$ input steps with a $H = 1$
output horizon across all experiments.

\subsection{Baseline Loss Functions}

We compare \RNAD against four established baselines:
\begin{itemize}
\item \textbf{MSE:} Standard mean-squared error; the canonical
  homoskedastic baseline against which all improvements are measured.
\item \textbf{Huber ($\delta = 1.0$):} Fixed-threshold interpolation
  between L2 and L1~\citep{huber1964robust}.
\item \textbf{log-cosh:} Smooth approximation to MAE with bounded
  second derivative~\citep{saleh2024logcosh}.
\item \textbf{GMADL ($\alpha = 0.5$, $\beta = 5.0$):} Combined
  magnitude and directional loss without volatility normalization
  or noise gating~\citep{michankow2024gmadl}.
\end{itemize}

\subsection{Forecasting Backbones}

The four evaluation backbones are:
\begin{itemize}
\item \textbf{PatchTST}~\citep{patchtst}: Patch-based transformer with
  channel-independent self-attention and RevIN normalization;
  $L = 24$, patch length $16$, stride $8$ (two patches),
  $d_{\mathrm{model}} = 128$, $d_{\mathrm{ff}} = 256$,
  $4$ heads, $3$ encoder layers, dropout $0.05$,
  fc-dropout $0.05$, and head-dropout $0.0$.
\item \textbf{iTransformer}~\citep{itransformer}: Inverted transformer
  with a 24-step linear inverted embedding, $d_{\mathrm{model}} = 128$,
  $d_{\mathrm{ff}} = 256$, $8$ heads, $3$ encoder layers, dropout
  $0.05$, and GELU activation.
\item \textbf{N-HiTS}~\citep{challu2023nhits}: Neural hierarchical
  interpolation with $3$ stacks, $2$ blocks per stack, $2$ MLP layers
  per block, $n_{\mathrm{hidden}} = 256$, max pooling, kernel sizes
  $[2,2,2]$, downsampling factors $[4,2,1]$, linear interpolation,
  dropout $0.05$, ReLU activation, no batch normalization, and no
  shared weights.
\item \textbf{TimeXer}~\citep{timexer2024}: Patch-based transformer with patch length
  $8$ (three non-overlapping target patches over $L=24$),
  $d_{\mathrm{model}} = 128$, $d_{\mathrm{ff}} = 256$,
  $8$ heads, $2$ encoder layers, dropout $0.05$,
  factor $5$, GELU activation, and \texttt{use\_norm=true}.
\end{itemize}
All backbones forecast the close channel (index 3) from the same
13-feature input vector. These architectures span transformer,
inverted-attention, hierarchical interpolation, and patch-based
cross-attention paradigms, providing a rigorous test of loss-level
generalization across diverse inductive biases.

\section{Implementation Details and Reproducibility}
\label{sec:reproducibility}

To ensure full reproducibility of the benchmark, we fix the exact
backbone, loss, optimizer, and deterministic settings used in the
codebase. The feature order and target index are identical across
all runs, and the same configuration file drives every experiment.

\subsection{Backbone Configurations}

All backbones forecast the close channel (index 3) from the same
24-step window of 13 standardized features. The exact settings are
summarized in Table~\ref{tab:architecture-config}.

\begin{table}[t]
\centering
\caption{Backbone configurations used in all benchmark runs.}
\label{tab:architecture-config}
\small
\begin{tabularx}{\textwidth}{l>{\raggedright\arraybackslash}X>{\raggedright\arraybackslash}X}
\toprule
Backbone & Exact configuration & Notes \\
\midrule
PatchTST & 3 encoder layers; 4 heads; $d_{\mathrm{model}} = 128$;
  $d_{\mathrm{ff}} = 256$; patch length 16; stride 8; dropout 0.05;
  fc-dropout 0.05; head-dropout 0.0 & RevIN on; decomposition off;
  BatchNorm backbone; GELU; \texttt{individual=false};
  \texttt{res\_attention=true}; \texttt{pe=zeros};
  \texttt{learn\_pe=true} \\
iTransformer & 3 encoder layers; 8 heads; $d_{\mathrm{model}} = 128$;
  $d_{\mathrm{ff}} = 256$; dropout 0.05; GELU activation & Inverted
  linear embedding over the 24-step window; LayerNorm encoder;
  attention factor 5; target index 3 \\
N-HiTS & 3 stacks; 2 blocks per stack; 2 MLP layers per block;
  $n_{\mathrm{hidden}} = 256$; dropout 0.05 & Max pooling; kernel
  sizes $[2,2,2]$; downsampling $[4,2,1]$; linear interpolation;
  ReLU; batch norm off; shared weights off \\
TimeXer & 2 encoder layers; 8 heads; $d_{\mathrm{model}} = 128$;
  $d_{\mathrm{ff}} = 256$; patch length 8; dropout 0.05; factor 5 &
  Three non-overlapping target patches over $L = 24$; global token;
  \texttt{use\_norm=true}; GELU \\
\bottomrule
\end{tabularx}
\end{table}

Unless otherwise noted, remaining constructor defaults follow the
model implementations in \texttt{models/*.py}. For example, PatchTST
uses \texttt{attn\_dropout=0}, \texttt{padding\_patch=None},
\texttt{affine=true}, \texttt{subtract\_last=false},
\texttt{pre\_norm=false}, \texttt{head\_type=flatten}, and
\texttt{max\_seq\_len=1024}.

\subsection{Optimization and Deterministic Training}

All models are trained with Adam~\citep{kingma2015adam} using the
PyTorch defaults ($\beta_1 = 0.9$, $\beta_2 = 0.999$,
$\epsilon = 10^{-8}$, zero weight decay, no AMSGrad) and a fixed
learning rate of $5 \times 10^{-4}$. No learning-rate schedule,
warmup, or gradient clipping is applied. Training runs for 50 epochs
with batch size $B = 64$ and \texttt{num\_workers = 0}. The training
split is shuffled; validation and test splits are not. Checkpoints
are saved every epoch, \texttt{resume=true} is enabled, and the best
checkpoint by validation RMSE is restored at the end of training.

\begin{table}[t]
\centering
\caption{Training and reproducibility settings shared across the
  benchmark.}
\label{tab:training-config}
\small
\begin{tabularx}{\textwidth}{l>{\raggedright\arraybackslash}X>{\raggedright\arraybackslash}X}
\toprule
Setting & Value & Notes \\
\midrule
Optimizer & Adam & PyTorch defaults: $\beta_1 = 0.9$, $\beta_2 = 0.999$,
  $\epsilon = 10^{-8}$, weight decay $= 0$, AMSGrad disabled \\
Learning rate & $5 \times 10^{-4}$ & Constant across all runs; no
  scheduler, warmup, or gradient clipping \\
Epochs & 50 & Full training budget; best validation-RMSE checkpoint
  restored after training \\
Batch size & 64 & Training loader shuffled; validation and test
  loaders unshuffled \\
Checkpointing & every epoch & \texttt{checkpoint\_every=1};
  \texttt{resume=true} \\
Seeds & $\{178, 956, 1045\}$ & Main benchmark uses three seeds;
  ablation studies use seed 178 \\
Determinism & true & Python, NumPy, PyTorch, and CUDA seeds are
  fixed; cuDNN deterministic paths are enabled \\
\bottomrule
\end{tabularx}
\end{table}

\subsection{RNAD Loss Hyperparameters}

\RNAD uses fixed scalar controls across all runs; the loss contains
no per-asset calibration or learned parameters. The EMA volatility
buffer is initialized to 1.0, updated online, and serialized in
checkpoints. Table~\ref{tab:hyperparams} lists the loss-specific
constants.

\begin{table}[t]
\centering
\caption{\RNAD loss hyperparameter configuration. All values are
  fixed across assets, backbones, and seeds. The noise gate is a
  smooth sigmoid rather than a hard threshold.}
\label{tab:hyperparams}
\small
\begin{tabularx}{\textwidth}{X c c c}
\toprule
Parameter                    & Value    & Role \\
\midrule
$\lambda_{\mathrm{dir}}$     & $0.5$    & Directional penalty weight (Eq.~\ref{eq:rnad-sample}) \\
$\tau_{\mathrm{mag}}$        & $1.1$    & Tail-gate transition midpoint (Eq.~\ref{eq:tail-gate}) \\
$\tau_{\mathrm{temp}}$       & $0.3$    & Tail-gate transition sharpness (Eq.~\ref{eq:tail-gate}) \\
$\varepsilon_c$              & $10^{-3}$ & Charbonnier smoothness offset (Eq.~\ref{eq:charb-kernel}) \\
$\beta$                      & $4.5$    & Tanh sharpness for sign encoding (Eq.~\ref{eq:tanh-sign}) \\
$\kappa$                     & $4.5$    & Softplus margin slope (Eq.~\ref{eq:dir-loss}) \\
$\tau_{\mathrm{noise}}$      & $0.10$   & Smooth noise-gate center (Eq.~\ref{eq:noise-gate}) \\
$\tau_{\mathrm{noise,temp}}$ & $0.04$   & Noise-gate temperature (Eq.~\ref{eq:noise-gate}) \\
$\rho$                       & $0.985$  & EMA momentum, $\approx 67$-step window (Eq.~\ref{eq:ema-scale}) \\
$\varepsilon$                & $10^{-8}$ & Numerical stability constant \\
\bottomrule
\end{tabularx}
\end{table}

\FloatBarrier

\subsection{Evaluation Metrics}

Predictions are inverse-transformed to raw close prices using the
fitted target scaler before RMSE, MAE, and directional accuracy (DA)
are computed.

\paragraph{Primary --- RMSE.}
Root Mean Squared Error on the raw price scale (lower is better):
\begin{equation}
  \mathrm{RMSE}
  = \sqrt{\frac{1}{|\mathcal{T}|}\sum_{t \in \mathcal{T}}
          \bigl(\hat{y}_t - y_t\bigr)^2}.
  \label{eq:rmse}
\end{equation}

\paragraph{Primary --- Directional Accuracy (DA).}
The proportion of non-flat steps for which the predicted return sign
matches the realized sign (higher is better):
\begin{equation}
  \mathrm{DA}
  = \frac{1}{|\mathcal{T}|}\sum_{t \in \mathcal{T}}
    \mathbf{1}\bigl[\sgn(\hat{r}_t) = \sgn(r_t)\bigr],
  \qquad
  \mathcal{T} = \bigl\{t : r_t \neq 0\bigr\}.
  \label{eq:da}
\end{equation}

\paragraph{Secondary.}
Mean Absolute Error (MAE) and ordinal RMSE rank across losses within
each backbone are reported as secondary diagnostics to assess
consistency.

All results are expressed as mean $\pm$ standard deviation over the
three benchmark seeds.

\FloatBarrier
% -----------------------------------------------------------------------
\section{Results}
\label{sec:results}

\subsection{Main Benchmark: Backbone--Loss Comparison}

Table~\ref{tab:main-bench} compares the five loss functions on all
four backbones, averaging over nine assets. \RNAD achieves the lowest
mean RMSE on all four backbones. Because RMSE is scale-dependent,
these cross-asset averages are influenced more by higher-priced assets;
the per-asset results in Table~\ref{tab:per-asset} should therefore
be read alongside the aggregate figures.
The mean RMSE rank of \RNAD is 1.0 on PatchTST, iTransformer, and
TimeXer, and 2.0 on N-HiTS. The $\Delta$RMSE column quantifies
improvement over the best non-\RNAD baseline; reductions range from
$-0.66$ (TimeXer) to $-0.88$ (PatchTST) on the three
transformer-family backbones. On N-HiTS, gains are positive but
smaller ($\Delta\mathrm{RMSE} = -0.086$), a pattern analyzed in
Section~\ref{ssec:interactions}.

Notably, \RNAD simultaneously improves directional accuracy relative
to MSE, Huber, log-cosh, and GMADL on all four backbones while also
reducing RMSE. This joint improvement suggests that directional and
magnitude objectives are not in tension under the gating design of
\RNAD; the noise gate appears to prevent directional supervision from
introducing gradient noise that would otherwise degrade magnitude
accuracy.

\begin{table}[!http]
\centering
\caption{Main benchmark results averaged over nine assets and three
  seeds ($\{178, 956, 1045\}$). Lower RMSE and MAE are better; higher
  DA~(\%) is better. Rank denotes ordinal position by RMSE within each
  backbone (1 = best). $\Delta$RMSE is the improvement of \RNAD over
  the best non-\RNAD baseline; negative values indicate improvement.
  Bold entries indicate the best-performing loss per backbone per
  metric.}
\label{tab:main-bench}
\scriptsize
\setlength{\tabcolsep}{3.0pt}
\resizebox{\textwidth}{!}{%
\begin{tabular}{llrrrrc}
\toprule
Model & Loss & RMSE & MAE & DA~(\%) & Rank & $\Delta$RMSE \\
\midrule
\multirow{5}{*}{PatchTST}
  & MSE      & $60.78\pm0.02$ & $40.36\pm0.05$ & $50.05\pm0.20$ & 3 & $+0.004$ \\
  & Huber    & $60.78\pm0.02$ & $40.36\pm0.05$ & $50.02\pm0.23$ & 4 & $+0.004$ \\
  & log-cosh & $\underline{60.77\pm0.02}$ & $40.35\pm0.04$ & $50.09\pm0.19$ & 2 & $0.000$  \\
  & GMADL    & $60.82\pm0.14$ & $\underline{40.27\pm0.17}$ & $\underline{50.33\pm0.01}$ & 5 & $+0.051$ \\
  & \RNAD    & $\mathbf{59.89\pm0.03}$ & $\mathbf{39.35\pm0.05}$
             & $\mathbf{50.97\pm0.14}$ & \textbf{1} & $\mathbf{-0.882}$ \\
\midrule
\multirow{5}{*}{iTransformer}
  & MSE      & $61.22\pm0.32$ & $40.62\pm0.15$ & $50.28\pm0.39$ & 3 & $+0.003$ \\
  & Huber    & $61.23\pm0.31$ & $40.61\pm0.19$ & $50.32\pm0.36$ & 4 & $+0.005$ \\
  & log-cosh & $\underline{61.22\pm0.31}$ & $\underline{40.61\pm0.18}$ & $50.22\pm0.06$ & 2 & $0.000$  \\
  & GMADL    & $61.67\pm0.41$ & $40.82\pm0.12$ & $\underline{50.51\pm0.49}$ & 5 & $+0.445$ \\
  & \RNAD    & $\mathbf{60.54\pm0.07}$ & $\mathbf{39.84\pm0.02}$
             & $\mathbf{51.29\pm0.10}$ & \textbf{1} & $\mathbf{-0.681}$ \\
\midrule
\multirow{5}{*}{N-HiTS}
  & MSE      & $\underline{60.20\pm0.05}$ & $39.90\pm0.13$ & $49.69\pm0.38$ & 2 & $0.000$ \\
  & Huber    & $60.47\pm0.38$ & $\underline{39.85\pm0.40}$
             & $\underline{50.06\pm0.39}$ & 4 & $+0.272$ \\
  & log-cosh & $60.40\pm0.29$ & $39.89\pm0.36$ & $49.83\pm0.55$ & 3 & $+0.202$ \\
  & GMADL    & $64.45\pm4.77$ & $42.39\pm2.84$ & $49.88\pm0.37$ & 5 & $+4.256$ \\
  & \RNAD    & $\mathbf{60.11\pm0.48}$ & $\mathbf{39.82\pm0.42}$ & $\mathbf{51.35\pm0.30}$ & \textbf{1} & $\mathbf{-0.086}$ \\
\midrule
\multirow{5}{*}{TimeXer}
  & MSE      & $60.80\pm0.18$ & $40.21\pm0.20$ & $50.23\pm0.11$ & 4 & $+0.029$ \\
  & Huber    & $60.78\pm0.08$ & $40.20\pm0.13$ & $50.19\pm0.26$ & 3 & $+0.009$ \\
  & log-cosh & $60.84\pm0.11$ & $40.22\pm0.14$ & $\underline{50.38\pm0.12}$ & 5 & $+0.069$ \\
  & GMADL    & $\underline{60.77\pm0.05}$ & $\underline{40.19\pm0.08}$ & $50.36\pm0.07$ & 2 & $0.000$ \\
  & \RNAD    & $\mathbf{60.12\pm0.01}$ & $\mathbf{39.58\pm0.02}$
             & $\mathbf{50.81\pm0.24}$ & \textbf{1} & $\mathbf{-0.658}$ \\
\bottomrule
\end{tabular}
}
\end{table}

\FloatBarrier

\begin{figure}[t]
\centering
\includegraphics[width=0.95\textwidth]{figure_02_rmse_vs_da_pareto_scatter.pdf}
\caption{\textbf{RMSE--DA Pareto frontier.}
  Each point represents one backbone--loss combination ($n = 20$
  combinations; RMSE and DA averaged across three seeds and nine
  assets). \RNAD (red) occupies the lower-right quadrant, achieving
  simultaneously lower RMSE and higher DA than MSE, Huber, and
  log-cosh baselines. GMADL improves DA at the cost of elevated
  RMSE on several backbones. RMSE is in raw price units (averaged
  over nine assets for visual comparison); DA in percent.}
\label{fig:pareto-frontier}
\end{figure}

\FloatBarrier

\subsection{Cross-Architecture Consistency}

Table~\ref{tab:model-wise} aggregates metrics by backbone across all
nine assets and three seeds. \RNAD achieves the best average RMSE on
all four architectures and reduces RMSE relative to MSE by
0.09--0.89 units. Standard deviations of per-seed RMSE on
transformer-family backbones remain small ($\sigma \leq 0.07$;
Table~\ref{tab:main-bench}), indicating that the observed gains are
stable across random seeds and are not attributable to favorable
initialization.

The contrast between transformer-family backbones and N-HiTS suggests
that backbone-level regime-aware inductive biases can partially address
the same failure mode targeted by \RNAD at the loss level. This
pattern indicates that loss-level and architecture-level
regime-awareness occupy partially substitutable functional roles, and
that \RNAD is likely most beneficial when paired with an expressive
but geometrically naive backbone.

\begin{table}[!htbp]
\centering
\caption{Backbone-aggregated performance (mean over nine assets and three seeds). $\Delta$ vs.\ MSE and $\Delta$ vs.\ GMADL are the RMSE improvements of \RNAD.}
\label{tab:model-wise}
\small
\resizebox{\textwidth}{!}{%
\begin{tabular}{llrrrr}
\toprule
Model        & Avg RMSE & Avg DA & $\Delta$ vs.\ MSE & $\Delta$ vs.\ GMADL & Best Loss \\
\midrule
PatchTST     & $\mathbf{60.61}$ & 50.29 & $\mathbf{-0.886}$ & $-0.933$ & \RNAD \\
iTransformer & 61.18 & $\mathbf{50.52}$ & $-0.684$ & $\underline{-1.126}$ & \RNAD \\
N-HiTS       & 61.08 & 50.16 & $-0.086$ & $\mathbf{-4.131}$ & \RNAD \\
TimeXer      & $\underline{60.66}$ & $\underline{50.39}$ & $\underline{-0.687}$ & $-0.658$ & \RNAD \\
\bottomrule
\end{tabular}
}
\end{table}

\FloatBarrier

\begin{figure}[!htbp]
\centering
\includegraphics[width=0.95\textwidth]{figure_03_cross_model_consistency.pdf}
\caption{\textbf{Cross-architecture consistency.}
  Average RMSE (bars, left axis) and directional accuracy (line
  markers, right axis) aggregated over nine assets and three seeds,
  grouped by backbone. \RNAD (red) reduces RMSE consistently on
  transformer-family backbones (PatchTST, iTransformer, TimeXer)
  relative to MSE, Huber, and log-cosh. On N-HiTS, the
  architecture's existing regime-aware inductive biases reduce the
  marginal benefit of loss-level intervention. Error bars denote
  standard deviation over seeds.}
\label{fig:cross-model}
\end{figure}

\FloatBarrier

\subsection{Per-Asset Analysis}

Table~\ref{tab:per-asset} reports per-asset RMSE and DA on PatchTST.
\RNAD achieves the lowest RMSE on all nine assets, indicating that
the gains are not concentrated within a single market class. RMSE
reductions are proportionally largest for BTCUSD ($-6.69$ units,
$-1.56\%$) and AAPL ($-0.133$ units, $-7.6\%$), the two assets with
the highest realized volatility (Table~\ref{tab:data-summary}),
consistent with the normalization mechanism providing the greatest
benefit in volatile regimes.

Directional accuracy improvements are more heterogeneous, ranging from
$+0.17$ percentage points (USDJPY) to $+2.29$ pp (BTCUSD). This
heterogeneity reflects both the inherent difficulty of one-step
directional prediction at the minute horizon and the asset-specific
relationship between volatility regime and the efficacy of gated
directional supervision.

\begin{table}[t]
\centering
\caption{Per-asset results on PatchTST (mean over seeds
  $\{178, 956, 1045\}$). RMSE (lower is better) and DA~(\%) (higher
  is better). Bold indicates best RMSE; underline indicates second
  best.}
\label{tab:per-asset}
\small
\setlength{\tabcolsep}{3pt}
\resizebox{\textwidth}{!}{%
\begin{tabular}{lcc cc cc cc cc}
\toprule
& \multicolumn{2}{c}{MSE}
& \multicolumn{2}{c}{Huber}
& \multicolumn{2}{c}{log-cosh}
& \multicolumn{2}{c}{GMADL}
& \multicolumn{2}{c}{\RNAD} \\
\cmidrule(lr){2-3} \cmidrule(lr){4-5} \cmidrule(lr){6-7} \cmidrule(lr){8-9} \cmidrule(lr){10-11}
Asset
& RMSE & DA
& RMSE & DA
& RMSE & DA
& RMSE & DA
& RMSE & DA \\
\midrule
AAPL    & 1.75 & 48.62 & 1.73 & 48.43 & 1.71 & 49.25 & \underline{1.67} & \underline{50.14} & \textbf{1.61} & \textbf{50.75} \\
AUDUSD  & 7.97e{-4} & 49.61 & \underline{7.97e{-4}} & 49.68 & 7.97e{-4} & \underline{49.74} & 8.02e{-4} & 49.67 & \textbf{7.90e{-4}} & \textbf{50.56} \\
BTCUSD  & \underline{427.90} & 50.01 & 427.92 & 49.98 & 427.90 & 50.00 & 428.12 & \underline{51.25} & \textbf{421.21} & \textbf{52.30} \\
DAX30   & 54.93 & \underline{50.62} & 54.93 & 50.62 & 54.93 & 50.59 & \underline{54.90} & 50.17 & \textbf{54.50} & \textbf{51.58} \\
ETHUSD  & 24.04 & 50.29 & \underline{24.03} & 50.21 & 24.03 & 50.28 & 24.25 & \underline{50.48} & \textbf{23.75} & \textbf{51.08} \\
FTSE100 & \underline{15.26} & 50.68 & 15.26 & 50.66 & 15.27 & 50.51 & 15.34 & \underline{50.86} & \textbf{15.06} & \textbf{51.17} \\
US500   & 14.47 & 49.98 & 14.47 & 49.85 & \underline{14.47} & 49.93 & 14.49 & \textbf{50.29} & \textbf{14.29} & \underline{50.26} \\
USDJPY  & 0.18 & 50.73 & 0.18 & 50.51 & \underline{0.18} & \underline{50.76} & 0.18 & 50.28 & \textbf{0.18} & \textbf{50.89} \\
XAUUSD  & \underline{8.47} & 49.93 & 8.47 & \textbf{50.23} & 8.47 & 49.73 & 8.47 & 49.86 & \textbf{8.41} & \underline{50.11} \\
\bottomrule
\end{tabular}%
}
\end{table}
\FloatBarrier

\begin{figure}[!htbp]
\centering
\includegraphics[width=0.95\textwidth]{figure_06_per_asset_metric_bars.pdf}
\caption{\textbf{Per-asset performance on PatchTST.}
  Left: RMSE by asset and loss; \RNAD (red) achieves the lowest
  RMSE in every asset class. Right: directional accuracy (DA, higher
  is better). RMSE axis scales differ substantially across assets
  due to absolute price-level variation (\eg BTCUSD vs.\ AUDUSD);
  percentage improvements are scale-consistent after volatility
  normalization.}
\label{fig:per-asset}
\end{figure}

\FloatBarrier

\subsection{Volatility Tracking Diagnostics}
\label{ssec:vol-tracking}

Figure~\ref{fig:regime-scale} shows the adaptive scale $s_t$
alongside $s_t^{\mathrm{batch}}$, $s_t^{\mathrm{EMA}}$, and a proxy
for realized volatility on BTCUSD and AUDUSD. The batch scale
responds immediately to volatility spikes; the EMA scale captures
regime-level persistence. Their geometric blend tracks realized
volatility accurately without overfitting to transient noise,
confirming that the normalization mechanism operates as intended
across assets with two-orders-of-magnitude differences in absolute
return scale.

\begin{figure}[!htbp]
\centering
\includegraphics[width=0.95\textwidth]{figure_05_regime_scale_tracking.pdf}
\caption{\textbf{Adaptive volatility tracking.}
  BTCUSD (top) and AUDUSD (bottom). Batch scale $s_t^{\mathrm{batch}}$
  (light blue, Eq.~\ref{eq:batch-scale}) reacts immediately to
  volatility spikes. EMA scale $s_t^{\mathrm{EMA}}$ (dark blue,
  Eq.~\ref{eq:ema-scale}, $\rho = 0.985$) provides regime-level
  memory. Composite scale $s_t$ (red, Eq.~\ref{eq:composite-scale}).
  Circular markers: realized volatility (rolling RMS return). The
  composite scale faithfully tracks realized volatility while
  attenuating transient noise.}
\label{fig:regime-scale}
\end{figure}

\FloatBarrier
% -----------------------------------------------------------------------
\section{Ablation Analysis}
\label{sec:ablation}

To assess each component's individual contribution, we train PatchTST
with seed 178 under four conditions, removing one mechanism at a time
from the full \RNAD objective:
\begin{itemize}
\item \texttt{FULL}: Complete \RNAD loss (Eq.~\ref{eq:rnad-sample}).
\item \texttt{NO\_DIR}: Directional term removed
  ($\lambda_{\mathrm{dir}} = 0$).
\item \texttt{NO\_TAIL}: Gated magnitude replaced by pure quadratic
  ($g_t \equiv 0$ in Eq.~\ref{eq:mag-loss}).
\item \texttt{NO\_NOISE}: Noise gate disabled ($w_t \equiv 1$ in
  Eq.~\ref{eq:noise-gate}; directional term always active).
\end{itemize}

Table~\ref{tab:ablation} reports results on four representative
assets: a high-volatility cryptocurrency (BTCUSD), a low-volatility
foreign exchange pair (AUDUSD), an equity index (US500), and a
commodity (XAUUSD).

\begin{table}[t]
\centering
\caption{Controlled ablation on PatchTST (seed 178). $\Delta$RMSE and
  $\Delta$DA are computed relative to \texttt{FULL} for each asset.
  Positive $\Delta$RMSE indicates degradation when a component is
  removed.}
\label{tab:ablation}
\small
\setlength{\tabcolsep}{2.5pt}
\resizebox{\textwidth}{!}{%
\begin{tabular}{llrrrrr}
\toprule
Asset  & Variant    & RMSE      & $\Delta$RMSE   & MAE      & DA    & $\Delta$DA \\
\midrule
\multirow{4}{*}{AUDUSD}
  & FULL       & $\underline{0.00079020}$   & ---            & $\underline{0.00053505}$  & 50.84 & ---    \\
  & NO\_DIR    & 0.00079655   & $+0.00000634$     & 0.00054031  & 50.24 & $-0.60$ \\
  & NO\_TAIL   & 0.00079142   & $+0.00000122$     & 0.00053625  & $\underline{50.86}$ & $+0.01$ \\
  & NO\_NOISE  & $\mathbf{0.00078971}$   & $-0.00000049$     & $\mathbf{0.00053444}$  & $\mathbf{51.15}$ & $+0.31$ \\
\midrule
\multirow{4}{*}{BTCUSD}
  & FULL       & 422.61    & ---            & $\underline{284.04}$   & 51.67 & ---    \\
  & NO\_DIR    & 425.32    & $+2.71$        & 286.09   & $\mathbf{51.74}$ & $+0.06$ \\
  & NO\_TAIL   & $\mathbf{421.87}$    & $-0.74$        & $\mathbf{283.54}$   & 51.51 & $-0.16$ \\
  & NO\_NOISE  & $\underline{422.42}$    & $-0.18$        & 284.27   & $\underline{51.72}$ & $+0.05$ \\
\midrule
\multirow{4}{*}{US500}
  & FULL       & $\mathbf{14.33}$     & ---            & $\underline{8.01}$     & 49.57 & ---    \\
  & NO\_DIR    & $\underline{14.33}$     & $+0.01$        & $\underline{8.06}$     & 48.60 & $-0.98$ \\
  & NO\_TAIL   & 14.35     & $+0.02$        & 8.03     & $\mathbf{50.27}$ & $+0.70$ \\
  & NO\_NOISE  & 14.34     & $+0.01$        & $\mathbf{8.00}$     & $\underline{49.97}$ & $+0.40$ \\
\midrule
\multirow{4}{*}{XAUUSD}
  & FULL       & $\underline{8.41}$      & ---            & 5.56     & $\underline{50.54}$ & ---    \\
  & NO\_DIR    & 8.48      & $+0.07$        & 5.60     & 50.43 & $-0.11$ \\
  & NO\_TAIL   & 8.42      & $+0.01$        & $\mathbf{5.56}$     & $\mathbf{50.90}$ & $+0.35$ \\
  & NO\_NOISE  & $\mathbf{8.41}$      & $-0.00$        & $\underline{5.56}$     & 50.39 & $-0.15$ \\
\bottomrule
\end{tabular}
}
\end{table}

\FloatBarrier

\paragraph{Directional learning (\texttt{NO\_DIR}).}
Removing the directional penalty consistently degrades RMSE (mean
$\Delta = +0.70$ RMSE units) and reduces DA on three of the four
assets. This indicates that directional supervision contributes to
both magnitude and directional accuracy---a coupling explained in
Section~\ref{ssec:coupling}. The effect is largest on BTCUSD
($\Delta = +2.71$), where economically significant directional moves
are most frequent and the noise gate therefore permits the most
directional supervision.

\paragraph{Tail gating (\texttt{NO\_TAIL}).}
The contribution of tail gating is asset-dependent. On BTCUSD,
removing the tail gate yields a marginal RMSE improvement ($-0.74$),
suggesting that $\tau_{\mathrm{mag}} = 1.1$ transitions slightly too
aggressively for BTCUSD's heavier tails relative to the default
calibration. On AUDUSD, US500, and XAUUSD, the full loss is preferred
over \texttt{NO\_TAIL}. These results indicate that
$\tau_{\mathrm{mag}}$ is the parameter most sensitive to
asset-specific tail behavior and is therefore a candidate for light
per-asset tuning in production deployments.

\paragraph{Noise gating (\texttt{NO\_NOISE}).}
The effects of noise gating are mixed and generally modest. On AUDUSD
and BTCUSD, disabling the gate slightly improves RMSE
($-4.9\times10^{-7}$ and $-0.18$, respectively) and increases DA.
On US500, the gate improves RMSE by $+0.01$ when removed while
reducing DA. On XAUUSD, disabling the gate yields a marginal RMSE
gain ($-0.001$) but reduces DA by $0.15$ pp. Overall, noise gating
functions as a conservative regularizer whose impact is
regime-dependent.

\begin{figure}[t]
\centering
\includegraphics[width=0.95\textwidth]{figure_04_ablation_comparison.pdf}
\caption{\textbf{Ablation component contributions.}
  Left: percent change in RMSE when each mechanism is removed from
  the full \RNAD objective, across four representative assets.
  Right: corresponding percent change in DA. \texttt{NO\_DIR}
  consistently degrades RMSE, indicating that directional supervision
  is the primary driver of magnitude accuracy improvements.
  \texttt{NO\_TAIL} is most impactful on heavy-tailed assets
  (BTCUSD, XAUUSD). \texttt{NO\_NOISE} exhibits mixed,
  low-magnitude effects across assets.}
\label{fig:ablation}
\end{figure}

\FloatBarrier
% -----------------------------------------------------------------------
\section{Discussion}
\label{sec:discussion}

\subsection{Why \RNAD Improves Forecasting}

The consistent RMSE improvements on transformer-style backbones are
interpretable through the lens of gradient signal quality. Standard
MSE generates parameter updates proportional to raw residual
magnitude, conflating prediction error with volatility scale. In
high-volatility regimes, large shocks dominate gradient updates,
driving the optimizer toward minimizing absolute magnitudes rather
than learning regime-relative structure. In low-volatility periods,
directional noise corrupts the sign gradient. \RNAD corrects both
pathologies simultaneously: volatility normalization calibrates
gradient magnitude to predictive difficulty; tail gating bounds the
influence of tail events; and gated directional supervision injects
complementary, sign-level information that MSE cannot express.

This framing is consistent with a signal-to-noise improvement in
the gradient stream. By normalizing residuals by $s_t$, the loss
ensures that gradient magnitudes reflect how unexpected a residual
is \emph{relative to the prevailing regime}, rather than how large
it is in absolute terms. The directional term then contributes a
complementary loss surface that creates asymmetric penalties around
the sign boundary, explicitly encoding the economically motivated
half-plane constraint that MSE ignores.

\subsection{Directional--Magnitude Coupling}
\label{ssec:coupling}

A notable empirical finding is that the addition of a directional
penalty also reduces RMSE (Table~\ref{tab:ablation}, \texttt{NO\_DIR}
rows). This coupling arises from the geometry of the squared-error
loss. For a prediction $\hat{r}_t$ with the wrong sign, the residual
is $e_t = \hat{y}_t - y_t$; when $|\hat{r}_t| \approx |r_t|$, a sign
error approximately doubles the residual compared to a same-magnitude
correct-sign prediction. Concretely, the average MSE contribution of
a sign-error event is roughly four times that of a correct-sign
prediction of equal magnitude. By penalizing sign errors, \RNAD
implicitly constrains $\hat{r}_t$ toward the correct half-plane,
reducing the frequency of large-residual sign-error events and thereby
improving MSE even when $\lambda_{\mathrm{dir}}$ is modest.

\subsection{Architectural Interactions}
\label{ssec:interactions}

The architecture-dependent benefit pattern---large gains on
transformer-family backbones, smaller gains on N-HiTS---supports the
interpretation that loss-level and architecture-level regime-awareness
occupy overlapping functional space. N-HiTS's hierarchical
decomposition already captures some regime structure; the marginal
benefit of \RNAD is correspondingly smaller when the backbone has
partially addressed the same failure mode. Importantly, \RNAD does
not degrade backbone-level aggregate RMSE: on N-HiTS the gain is
modest but positive ($\Delta\mathrm{RMSE} = -0.086$ versus MSE),
confirming that the loss is safely composable with regime-aware
architectures.

\subsection{Statistical Stability}

The small standard deviations of \RNAD RMSE across seeds
($\sigma \leq 0.07$ on transformer backbones; Table~\ref{tab:model-wise})
confirm that the observed improvements are stable across random
initializations. This stability is consistent with the gradient-level
explanation above: improvements arise from systematic changes in the
optimization geometry rather than from favorable sampling of the
loss landscape.

\FloatBarrier
% -----------------------------------------------------------------------
\section{Limitations}
\label{sec:limitations}

\paragraph{Theoretical foundations.}
\RNAD is designed on geometric and empirical grounds. No convergence
guarantees, Bayes-optimality certificates, or formal proof that the
gated directional term constitutes a proper scoring
rule~\citep{gneiting2007strictly} are provided. The EMA volatility
estimator introduces first-order autocorrelation in the effective loss
landscape and may misrepresent long-memory volatility dynamics not
captured by the GARCH-like EMA kernel. Formal Diebold--Mariano or
Model Confidence Set tests~\citep{diebold1995comparing,
hansen2011model} were not applied; extending these results with such
statistical comparisons is an important direction for future work.

\paragraph{Hyperparameter sensitivity.}
The loss introduces ten scalar controls (Table~\ref{tab:hyperparams}).
While the default values are empirically robust to $\pm 20\%$
perturbations for most parameters, $\tau_{\mathrm{noise}}$ is the
most sensitive control and may merit retuning when transferring to a
substantially different market universe. The tail-gating parameters
$\tau_{\mathrm{mag}}$ and $\tau_{\mathrm{temp}}$ are moderately
sensitive on heavy-tailed assets (see the BTCUSD ablation results
in Table~\ref{tab:ablation}) and represent the most valuable targets
for asset-specific calibration.

\paragraph{Scope of demonstrated improvements.}
Consistent gains are concentrated in one-step-ahead, raw-price
prediction on attention-based architectures. Benefits are weaker
or absent on regime-aware architectures, at longer forecasting
horizons, and in return-domain or log-price formulations where
normalization is less critical. Practitioners should validate \RNAD
on their specific setup before assuming transferability.

\paragraph{Economic validity of directional accuracy.}
The DA metric (Eq.~\ref{eq:da}) weights all non-flat moves equally
regardless of magnitude or economic significance, and does not account
for transaction costs, bid--ask spread, slippage, position sizing,
or the asymmetric utility of large versus small directional errors.
Higher DA does not imply higher risk-adjusted returns, and this study
makes no such claim. Evaluation within a full trading simulation with
realistic costs is necessary to draw conclusions about practical
financial value.

\paragraph{Statefulness and engineering overhead.}
The EMA buffer must be synchronized across data-parallel workers,
persisted at checkpoints, and reset at each training initialization.
This engineering overhead is absent in stateless losses and may
complicate deployment in highly distributed training pipelines.

\paragraph{Batch dependence.}
Because the volatility scale and gates are computed from mini-batches,
the objective is mildly batch-size dependent; smaller batches yield
noisier scale estimates and can alter the effective gate values.
Batch size is held fixed at 64 in all experiments, which controls
but does not eliminate this effect.

\paragraph{Evaluation breadth.}
The benchmark uses three seeds and a single chronological split per
asset. Extending to multiple non-overlapping temporal windows and
applying bootstrap inference over test-set samples would substantially
strengthen the statistical conclusions.

\FloatBarrier
% -----------------------------------------------------------------------
\section{Conclusion}
\label{sec:conclusion}

We introduced \RNAD, a regime-normalized adaptive directional loss
for financial time-series forecasting. By unifying adaptive volatility
normalization, tail-robust magnitude handling, and gated directional
learning into a single differentiable objective, \RNAD realigns the
training geometry with the distinctive statistical structure of
financial markets: regime-switching, heavy-tailed, and sign-sensitive.

Comprehensive benchmarks across nine assets and four architectures
demonstrate consistent RMSE improvements on transformer-style
backbones and the lowest per-asset RMSE on every PatchTST market.
Controlled ablation studies indicate that all three components
contribute independently: directional learning is the primary driver
of magnitude improvement; tail gating is most critical on
heavy-tailed assets; and noise gating contributes modest,
regime-dependent regularization. The loss is computationally
lightweight, architecture-agnostic, and requires no modification
to model structure, data pipeline, or inference code.

The broader implication is that \emph{loss-level design is an
important lever for financial forecasting improvement}, with effects
in this benchmark that are of a comparable order of magnitude to
several architectural choices while remaining less costly to deploy.
Directional accuracy should be interpreted as a forecast-quality
metric rather than a direct proxy for trading profitability; this
study does not evaluate transaction costs, slippage, or
profit-and-loss. When the forecasting target is volatile,
heavy-tailed, and direction-sensitive---as financial prices
characteristically are---the training objective can be highly
consequential relative to the choice of model architecture.

\FloatBarrier

% -----------------------------------------------------------------------
\appendix
\section{Supplementary Diagnostics}

\subsection{Loss Landscape and Gradient Dynamics}

Figure~\ref{fig:gradient-field} contrasts the gradient vector fields
of \RNAD and MSE as functions of normalized error $u_t$ and realized
move $z_t$. \RNAD maintains steep, informative gradients near the
origin (facilitating precise level learning) while transitioning to
sub-linear damping in the tails (bounding gradient norms during
extreme events). MSE gradients grow linearly without bound,
systematically over-weighting tail observations in parameter updates.

\begin{figure}[!htbp]
\centering
\includegraphics[width=0.95\textwidth]{figure_07_loss_landscape_gradient_field.pdf}
\caption{\textbf{Loss landscape and gradient fields: \RNAD vs.\ MSE.}
  (A) \RNAD maintains steep gradients for $|u_t| < 1$ and transitions
  to sub-linear damping for $|u_t| > 2$. (B) MSE gradients grow
  proportionally to $|u_t|$ without bound. Arrows indicate direction
  and relative magnitude of local optimization steps; color indicates
  loss density. Transition midpoint $\tau_{\mathrm{mag}} = 1.1$.}
\label{fig:gradient-field}
\end{figure}

\FloatBarrier

\subsection{Regime-Stratified Performance}

Figure~\ref{fig:volatility-regimes} stratifies test performance by
causal volatility regime (terciles of realized volatility computed
from training-set statistics, to avoid look-ahead bias). \RNAD
maintains a consistent performance advantage over MSE across all three
regimes, indicating that adaptive normalization generalizes across
volatility levels and does not merely exploit tail events.

\begin{figure}[!htbp]
\centering
\includegraphics[width=0.95\textwidth]{figure_08_performance_vs_volatility_regime.pdf}
\caption{\textbf{Performance stratified by volatility regime.}
  RMSE (bars, left axis) and DA (line, right axis) on BTCUSD, US500,
  and XAUUSD, stratified into low- ($<$33rd pctile), medium-
  (33rd--66th pctile), and high-volatility ($>$66th pctile) terciles.
  \RNAD (red) maintains a consistent advantage over MSE in all bins,
  supporting the regime-uniform generalization of the EMA
  normalization.}
\label{fig:volatility-regimes}
\end{figure}

\FloatBarrier

\subsection{Seed Stability}

Figure~\ref{fig:seed-stability} shows box-and-whisker distributions
of RMSE and DA across nine assets and three seeds. \RNAD distributions
are consistently shifted toward lower error, with narrow whiskers
indicating that improvements arise from the loss geometry rather than
favorable initialization.

\begin{figure}[!htbp]
\centering
\includegraphics[width=0.95\textwidth]{figure_09_stability_across_seeds.pdf}
\caption{\textbf{Seed stability.}
  Distributions of RMSE (left), MAE (center), and DA (right)
  for \RNAD vs.\ MSE, over $n = 9$ assets and three seeds.
  \RNAD distributions shift consistently toward lower error, with
  narrow inter-seed spread ($< 0.1\%$ relative variance on
  transformer backbones).}
\label{fig:seed-stability}
\end{figure}

\FloatBarrier

\subsection{Forecast Trace Case Studies}

Figure~\ref{fig:case-studies} shows 100-step forecast traces for
four assets. \RNAD tracks directional reversals more faithfully than
MSE during high-volatility windows ($|z_t| > \tau_{\mathrm{noise}}$),
consistent with the quantitative DA improvements reported in
Table~\ref{tab:per-asset}.

\begin{figure}[!htbp]
\centering
\includegraphics[width=0.95\textwidth]{figure_10_forecast_trace_case_studies.pdf}
\caption{\textbf{Forecast trace case studies.}
  PatchTST predictions (MSE: blue; \RNAD: red) vs.\ realized close
  price (black) over 100-step windows on four assets. Shaded
  regions highlight high-volatility episodes where \RNAD captures
  directional reversals that MSE misses. Right panels show 20-step
  rolling DA.}
\label{fig:case-studies}
\end{figure}

\FloatBarrier

\subsection{Full Ablation: All Assets and Components}

Figures~\ref{fig:variant-delta} and~\ref{fig:component-effects} extend
the four-asset ablation reported in Section~\ref{sec:ablation}.

\begin{figure}[!htbp]
\centering
\includegraphics[width=0.95\textwidth]{figure_A1_variant_delta_plot.pdf}
\caption{\textbf{Per-asset ablation variant delta.}
  Mean RMSE difference ($\Delta$RMSE vs.\ \texttt{FULL}) across the
  available ablation runs (seed 178) for the four reported assets.
  All per-asset mean deltas are non-negative, indicating that no
  individual component is redundant across the reported assets.}
\label{fig:variant-delta}
\end{figure}

\FloatBarrier

\begin{figure}[!htbp]
\centering
\includegraphics[width=0.95\textwidth]{figure_A2_component_effect_heatmap.pdf}
\caption{\textbf{Component contribution heatmap.}
  Mean percent RMSE increase when each mechanism is ablated, by asset
  and component (NORM: normalization; TAIL: tail gating; DIR:
  directional term; NOISE: noise gating). Darker cells indicate
  greater reliance on that component. The directional term (DIR)
  shows a strong and consistent contribution across assets; tail
  gating (TAIL) is most critical for high-volatility assets
  (BTCUSD, ETHUSD).}
\label{fig:component-effects}
\end{figure}

\FloatBarrier

\subsection{Computational Overhead}

\RNAD introduces negligible runtime cost relative to MSE. The EMA
update (Eq.~\ref{eq:ema-scale}) requires a single scalar
multiply-add per gradient step. Additional loss computations---tanh,
logistic gate, softplus---are elementwise operations parallelized by
the autograd engine with no synchronization barrier. Memory overhead
amounts to one floating-point scalar per asset. In wall-clock
benchmarks on all four backbones, training and inference throughput
are indistinguishable from MSE at standard batch sizes.

% -----------------------------------------------------------------------
\section*{Acknowledgments}

The author have no specific acknowledgments to report. No external
funding was received for this work, and no individuals or organizations
contributed in a manner that warrants formal acknowledgment.

\section*{Conflict of Interest}
The author declares that there are no conflicts of interest, financial or otherwise, that could have influenced the research, authorship, or publication of this work. The study was conducted independently, and no external funding, affiliations, or personal relationships exist that could be perceived as affecting the objectivity, integrity, or interpretation of the results presented in this article.

\section*{Data and Code Availability}

The implementation code, along with all relevant configuration files required to reproduce the experiments presented in this manuscript, is publicly available at the following GitHub repository:

\url{https://github.com/NabeelAhmad9/RNAD}

Processed data artifacts generated and analyzed during this study are not publicly available due to proprietary or ethical restrictions. They are available from the corresponding author upon reasonable request, subject to a signed data use agreement.
% -----------------------------------------------------------------------
\bibliographystyle{plainnat}
\begin{thebibliography}{99}

% --- Original references (revised/corrected) ---

\bibitem[Barron(2019)]{barron2019general}
J.~T. Barron.
A general and adaptive robust loss function.
In \emph{Proc.\ IEEE/CVF CVPR}, pages 4331--4339, 2019.

\bibitem[Bengio et~al.(2009)]{bengio2009curriculum}
Y.~Bengio, J.~Louradour, R.~Collobert, and J.~Weston.
Curriculum learning.
In \emph{Proc.\ ICML}, pages 41--48, 2009.

\bibitem[Bollerslev(1986)]{bollerslev1986generalized}
T.~Bollerslev.
Generalized autoregressive conditional heteroskedasticity.
\emph{Journal of Econometrics}, 31(3):307--327, 1986.

\bibitem[Cao and Tay(2001)]{cao2001financial}
L.~J. Cao and F.~E.~H. Tay.
Financial forecasting using support vector machines.
\emph{Neural Computing \& Applications}, 10(2):184--192, 2001.

\bibitem[Challu et~al.(2023)]{challu2023nhits}
C.~Challu, K.~G. Olivares, B.~N. Oreshkin, F.~Garza,
M.~Mergenthaler-Canseco, and A.~Dubrawski.
N-HiTS: Neural hierarchical interpolation for time series forecasting.
In \emph{Proc.\ AAAI}, volume~37, pages 6989--6997, 2023.

\bibitem[Charbonnier et~al.(1997)]{charbonnier1994deterministic}
P.~Charbonnier, L.~Blanc-F{\'e}raud, G.~Aubert, and M.~Barlaud.
Deterministic edge-preserving regularization in computed imaging.
\emph{IEEE Trans.\ Image Process.}, 6(2):298--311, 1997.

\bibitem[Saleh and Saleh(2024)]{saleh2024logcosh}
R.~A. Saleh and A.~K.~Md. Ehsanes Saleh.
Statistical Properties of the log-cosh Loss Function Used in Machine Learning.
\emph{arXiv preprint arXiv:2208.04564}, 2024.

\bibitem[Diebold and Mariano(1995)]{diebold1995comparing}
F.~X. Diebold and R.~S. Mariano.
Comparing predictive accuracy.
\emph{Journal of Business \& Economic Statistics}, 13(3):253--263, 1995.

\bibitem[Ekambaram et~al.(2023)]{ekambaram2023tsmixer}
V.~Ekambaram, A.~Jati, N.~Nguyen, P.~Sinthong, and J.~Kalagnanam.
TSMixer: Lightweight MLP-mixer model for multivariate time series forecasting.
In \emph{Proc.\ ACM KDD}, 2023.

\bibitem[Engle(1982)]{engle1982autoregressive}
R.~F. Engle.
Autoregressive conditional heteroskedasticity with estimates of the variance
of United Kingdom inflation.
\emph{Econometrica}, 50(4):987--1007, 1982.

\bibitem[Gneiting and Raftery(2007)]{gneiting2007strictly}
T.~Gneiting and A.~E. Raftery.
Strictly proper scoring rules, prediction, and estimation.
\emph{Journal of the American Statistical Association},
102(477):359--378, 2007.

\bibitem[Hansen et~al.(2011)]{hansen2011model}
P.~R. Hansen, A.~Lunde, and J.~M. Nason.
The model confidence set.
\emph{Econometrica}, 79(2):453--497, 2011.

\bibitem[Hastie et~al.(2009)]{hastie2009elements}
T.~Hastie, R.~Tibshirani, and J.~Friedman.
\emph{The Elements of Statistical Learning}, 2nd ed.
Springer, New York, 2009.

\bibitem[Huber(1964)]{huber1964robust}
P.~J. Huber.
Robust estimation of a location parameter.
\emph{Annals of Mathematical Statistics}, 35(1):73--101, 1964.

\bibitem[Kendall et~al.(2018)]{kendall2018multi}
A.~Kendall, Y.~Gal, and R.~Cipolla.
Multi-task learning using uncertainty to weigh losses for scene geometry
and semantics.
In \emph{Proc.\ IEEE/CVF CVPR}, pages 7482--7491, 2018.

\bibitem[Kim et~al.(2022)]{kim2022reversible}
T.~Kim, J.~Kim, Y.~Tae, C.~Park, J.-H. Choi, and J.~Choo.
Reversible instance normalization for accurate time-series forecasting
against distribution shift.
In \emph{Proc.\ ICLR}, 2022.

\bibitem[Kingma and Ba(2015)]{kingma2015adam}
D.~P. Kingma and J.~Ba.
Adam: A method for stochastic optimization.
In \emph{Proc.\ ICLR}, 2015.

\bibitem[Koenker and Bassett(1978)]{koenker1978regression}
R.~Koenker and G.~Bassett.
Regression quantiles.
\emph{Econometrica}, 46(1):33--50, 1978.

\bibitem[Kuan and White(1994)]{kuan1995forecasting}
C.-M. Kuan and H.~White.
Artificial neural networks: An econometric perspective.
\emph{Econometric Reviews}, 13(1):1--91, 1994.

\bibitem[Kumar et~al.(2010)]{kumar2010self}
M.~P. Kumar, B.~Packer, and D.~Koller.
Self-paced learning for latent variable models.
In \emph{Proc.\ NeurIPS}, pages 1189--1197, 2010.

\bibitem[Leitch and Tanner(1991)]{leitch1991economic}
G.~Leitch and J.~E. Tanner.
Economic forecast evaluation: Profits versus the conventional error measures.
\emph{American Economic Review}, 81(3):580--590, 1991.

\bibitem[Lim et~al.(2021)]{lim2021temporal}
B.~Lim, S.~O. Arik, N.~Loeff, and T.~Pfister.
Temporal fusion transformers for interpretable multi-horizon time series
forecasting.
\emph{International Journal of Forecasting}, 37(4):1748--1764, 2021.

\bibitem[Lin et~al.(2017)]{focal}
T.-Y. Lin, P.~Goyal, R.~Girshick, K.~He, and P.~Doll{\'a}r.
Focal loss for dense object detection.
In \emph{Proc.\ IEEE ICCV}, pages 2980--2988, 2017.

\bibitem[Liu et~al.(2024)]{itransformer}
Y.~Liu, T.~Hu, H.~Zhang, C.~Wu, S.~Wang, L.~Ma, and M.~Long.
iTransformer: Inverted transformers are effective for time series forecasting.
In \emph{Proc.\ ICLR}, 2024.

\bibitem[Wang et~al.(2024)]{timexer2024}
Y.~Wang, H.~Wu, J.~Dong, G.~Qin, H.~Zhang, Y.~Liu, Y.~Qiu, J.~Wang,
and M.~Long.
TimeXer: Empowering transformers for time series forecasting with exogenous
variables.
\emph{arXiv preprint arXiv:2402.19072}, 2024.

\bibitem[Makridakis et~al.(2022)]{makridakis2022m5}
S.~Makridakis, E.~Spiliotis, and V.~Assimakopoulos.
M5 accuracy competition: Results, findings, and conclusions.
\emph{International Journal of Forecasting}, 38(4):1346--1364, 2022.

\bibitem[Nie et~al.(2023)]{patchtst}
Y.~Nie, N.~H. Nguyen, P.~Sinthong, and J.~Kalagnanam.
A time series is worth 64 words: Long-term forecasting with transformers.
In \emph{Proc.\ ICLR}, 2023.

\bibitem[van~den Oord et~al.(2016)]{oord2016wavenet}
A.~van~den Oord, S.~Dieleman, H.~Zen, K.~Simonyan, O.~Vinyals, A.~Graves,
N.~Kalchbrenner, A.~Senior, and K.~Kavukcuoglu.
WaveNet: A generative model for raw audio.
\emph{arXiv:1609.03499}, 2016.

\bibitem[Oreshkin et~al.(2020)]{oreshkin2020nbeats}
B.~N. Oreshkin, D.~Carpov, N.~Chapados, and Y.~Bengio.
N-BEATS: Neural basis expansion analysis for interpretable time series
forecasting.
In \emph{Proc.\ ICLR}, 2020.

\bibitem[Pesaran and Timmermann(1992)]{pesaran1992simple}
M.~H. Pesaran and A.~Timmermann.
A simple nonparametric test of predictive performance.
\emph{Journal of Business \& Economic Statistics}, 10(4):461--465, 1992.

\bibitem[Salinas et~al.(2020)]{salinas2020deepar}
D.~Salinas, V.~Flunkert, J.~Gasthaus, and T.~Januschowski.
DeepAR: Probabilistic forecasting with autoregressive recurrent networks.
\emph{International Journal of Forecasting}, 36(3):1181--1191, 2020.

\bibitem[Sezer et~al.(2020)]{sezer2020financial}
O.~B. Sezer, M.~U. Gudelek, and A.~M. Ozbayoglu.
Financial time series forecasting with deep learning: A systematic
literature review.
\emph{Applied Soft Computing}, 90:106181, 2020.

\bibitem[Wang et~al.(2024)]{wang2024timemixer}
S.~Wang, H.~Wu, X.~Shi, T.~Hu, H.~Luo, L.~Ma, J.~Y. Zhang, and J.~Zhou.
TimeMixer: Decomposable multiscale mixing for time series forecasting.
In \emph{Proc.\ ICLR}, 2024.

\bibitem[White(1988)]{white1988economic}
H.~White.
Economic prediction using neural networks: The case of IBM daily stock
returns.
In \emph{Proc.\ IEEE ICNN}, volume~2, pages 451--458, 1988.

\bibitem[Wu et~al.(2021)]{wu2021autoformer}
H.~Wu, J.~Xu, J.~Wang, and M.~Long.
Autoformer: Decomposition transformers with auto-correlation for long-term
series forecasting.
In \emph{Proc.\ NeurIPS}, pages 22419--22430, 2021.

\bibitem[Wu et~al.(2023)]{wu2023timesnet}
H.~Wu, T.~Hu, Y.~Liu, H.~Zhou, J.~Wang, and M.~Long.
TimesNet: Temporal 2D-variation modeling for general time series analysis.
In \emph{Proc.\ ICLR}, 2023.

\bibitem[Micha{\'n}k{\'o}w et~al.(2024)]{michankow2024gmadl}
J.~Micha{\'n}k{\'o}w, P.~Sakowski, and R.~\'{S}lepaczuk.
Generalized Mean Absolute Directional Loss as a Solution to Overfitting and
High Transaction Costs in Machine Learning Models Used in High-Frequency
Algorithmic Investment Strategies.
\emph{arXiv preprint arXiv:2412.18405}, 2024.

\bibitem[Zeng et~al.(2023)]{zeng2023dlinear}
A.~Zeng, M.~Chen, L.~Zhang, and Q.~Xu.
Are transformers effective for time series forecasting?
In \emph{Proc.\ AAAI}, volume~37, pages 11121--11128, 2023.

\bibitem[Zhou et~al.(2021)]{zhou2021informer}
H.~Zhou, S.~Zhang, J.~Peng, S.~Zhang, J.~Li, H.~Xiong, and W.~Zhang.
Informer: Beyond efficient transformer for long sequence time-series
forecasting.
In \emph{Proc.\ AAAI}, volume~35, pages 11106--11115, 2021.

\bibitem[Zhou et~al.(2022)]{zhou2022fedformer}
T.~Zhou, Z.~Ma, Q.~Wen, X.~Wang, L.~Sun, and R.~Jin.
FEDformer: Frequency enhanced decomposed transformer for long-term series
forecasting.
In \emph{Proc.\ ICML}, pages 27268--27286, 2022.

% --- New references (25+) ---

\bibitem[Ansari et~al.(2024)]{ansari2024chronos}
A.~F. Ansari, L.~Stella, C.~Turkmen, X.~Zhang, P.~Mercado, H.~Shen,
O.~Shchur, S.~S. Rangapuram, S.~P. Pineda Arango, S.~Kapoor,
J.~Zschiegner, P.~Maddix, M.~W. Mahoney, K.~Januschowski, and
S.~Wang.
Chronos: Learning the language of time series.
\emph{arXiv:2403.07815}, 2024.

\bibitem[Burges et~al.(2005)]{burges2005learning}
C.~Burges, T.~Shaked, E.~Renshaw, A.~Lazier, M.~Deeds, N.~Hamilton,
and G.~Hullender.
Learning to rank using gradient descent.
In \emph{Proc.\ ICML}, pages 89--96, 2005.

\bibitem[Campbell and Thompson(2008)]{campbell2008predicting}
J.~Y. Campbell and S.~B. Thompson.
Predicting excess stock returns out of sample: Can anything beat the
historical average?
\emph{Review of Financial Studies}, 21(4):1509--1531, 2008.

\bibitem[Clark and West(2007)]{clark2007approximately}
T.~E. Clark and K.~D. West.
Approximately normal tests for equal predictive accuracy in nested models.
\emph{Journal of Econometrics}, 138(1):291--311, 2007.

\bibitem[Cont(2001)]{cont2001empirical}
R.~Cont.
Empirical properties of asset returns: Stylized facts and statistical
issues.
\emph{Quantitative Finance}, 1(2):223--236, 2001.

\bibitem[Das et~al.(2023)]{das2023long}
A.~Das, W.~Kong, A.~Sen, and Y.~Zhou.
Long-term forecasting with TiDE: Time-series dense encoder.
\emph{arXiv:2304.08424}, 2023.

\bibitem[Das et~al.(2024)]{das2024decoder}
A.~Das, W.~Kong, R.~Leach, S.~Mathur, A.~Sen, and R.~Yu.
Decoder-only foundation model for time-series forecasting.
In \emph{Proc.\ ICML}, 2024.

\bibitem[Diebold(2015)]{diebold2015comparing}
F.~X. Diebold.
Comparing predictive accuracy, twenty years later: A personal perspective on
the use and abuse of Diebold--Mariano tests.
\emph{Journal of Business \& Economic Statistics}, 33(1):1--24, 2015.

\bibitem[Fan et~al.(2023)]{fan2023dish}
W.~Fan, S.~Wang, Y.~Wang, X.~Wang, H.~Zhao, T.~Liu, C.~Shen, Q.~Xu,
G.~Liu, and X.~Wang.
Dish-TS: A general paradigm for alleviating distribution shift in time
series forecasting.
In \emph{Proc.\ AAAI}, volume~37, pages 7522--7529, 2023.

\bibitem[Freyberger et~al.(2020)]{freyberger2020dissecting}
J.~Freyberger, A.~Neuhierl, and M.~Weber.
Dissecting characteristics nonparametrically.
\emph{Review of Financial Studies}, 33(5):2326--2377, 2020.

\bibitem[Ganin et~al.(2016)]{ganin2016domain}
Y.~Ganin, E.~Ustinova, H.~Ajakan, P.~Germain, H.~Larochelle, F.~Laviolette,
M.~Marchand, and V.~Lempitsky.
Domain-adversarial training of neural networks.
\emph{Journal of Machine Learning Research}, 17(59):1--35, 2016.

\bibitem[Gretton et~al.(2012)]{gretton2012kernel}
A.~Gretton, K.~M. Borgwardt, M.~J. Rauch, B.~Sch{\"o}lkopf, and
A.~Smola.
A kernel two-sample test.
\emph{Journal of Machine Learning Research}, 13(25):723--773, 2012.

\bibitem[Gu et~al.(2020)]{gu2020empirical}
S.~Gu, B.~Kelly, and D.~Xiu.
Empirical asset pricing via machine learning.
\emph{Review of Financial Studies}, 33(5):2223--2273, 2020.

\bibitem[Guti{\'e}rrez et~al.(2016)]{gutierrez2016ordinal}
P.~A. Guti{\'e}rrez and M.~P{\'e}rez-Ortiz.
Ordinal regression methods: Survey and experimental study.
\emph{IEEE Trans.\ Knowledge and Data Engineering}, 28(1):127--146, 2016.

\bibitem[Hamilton(1989)]{hamilton1989new}
J.~D. Hamilton.
A new approach to the economic analysis of nonstationary time series and the
business cycle.
\emph{Econometrica}, 57(2):357--384, 1989.

\bibitem[Harvey et~al.(1997)]{harvey1997testing}
D.~I. Harvey, S.~J. Leybourne, and P.~Newbold.
Testing the equality of prediction mean squared errors.
\emph{International Journal of Forecasting}, 13(2):281--291, 1997.

\bibitem[Jin et~al.(2023)]{jin2023time}
M.~Jin, S.~Wang, L.~Ma, Z.~Chu, J.~Y. Zhang, X.~Shi, P.-Y. Chen,
Y.~Liang, Y.-F. Li, S.~Pan, and Q.~Wen.
Time-LLM: Time series forecasting by reprogramming large language models.
In \emph{Proc.\ ICLR}, 2024.

\bibitem[Kelly et~al.(2019)]{kelly2019characteristics}
B.~Kelly, S.~Pruitt, and Y.~Su.
Characteristics are covariances: A unified model of risk and return.
\emph{Journal of Financial Economics}, 134(3):501--524, 2019.

\bibitem[Liu et~al.(2022)]{liu2022non}
Y.~Liu, H.~Wu, J.~Wang, and M.~Long.
Non-stationary transformers: Exploring the stationarity in time series
forecasting.
In \emph{Proc.\ NeurIPS}, pages 9881--9893, 2022.

\bibitem[Liu et~al.(2024)]{liu2024adaptive}
Y.~Liu, X.~Hu, H.~Zhang, M.~Li, Z.~Guo, W.~Zhao, and M.~Long.
Adaptive normalization for non-stationary time series forecasting:
A temporal slice perspective.
In \emph{Proc.\ NeurIPS}, 2024.

\bibitem[Lo and MacKinlay(1988)]{lo1988stock}
A.~W. Lo and A.~C. MacKinlay.
Stock market prices do not follow random walks: Evidence from a simple
specification test.
\emph{Review of Financial Studies}, 1(1):41--66, 1988.

\bibitem[Mandelbrot(1963)]{mandelbrot1963variation}
B.~Mandelbrot.
The variation of certain speculative prices.
\emph{Journal of Business}, 36(4):394--419, 1963.

\bibitem[Meinshausen and Ridgeway(2006)]{meinshausen2006quantile}
N.~Meinshausen and G.~Ridgeway.
Quantile regression forests.
\emph{Journal of Machine Learning Research}, 7:983--999, 2006.

\bibitem[Ren et~al.(2018)]{ren2018learning}
M.~Ren, W.~Zeng, B.~Yang, and R.~Urtasun.
Learning to reweight examples for robust deep learning.
In \emph{Proc.\ ICML}, pages 4334--4343, 2018.

\bibitem[Stan{k}eviciute et~al.(2021)]{stankeviciute2021conformal}
K.~Stan{k}eviciute, A.~M. Alaa, and M.~van~der Schaar.
Conformal time-series forecasting.
In \emph{Proc.\ NeurIPS}, pages 6216--6228, 2021.

\bibitem[Vaswani et~al.(2017)]{vaswani2017attention}
A.~Vaswani, N.~Shazeer, N.~Parmar, J.~Uszkoreit, L.~Jones,
A.~N. Gomez, {\L}.~Kaiser, and I.~Polosukhin.
Attention is all you need.
In \emph{Proc.\ NeurIPS}, pages 5998--6008, 2017.

\bibitem[Welch and Goyal(2008)]{welch2008comprehensive}
I.~Welch and A.~Goyal.
A comprehensive look at the empirical performance of equity premium prediction.
\emph{Review of Financial Studies}, 21(4):1455--1508, 2008.

\bibitem[Woo et~al.(2022)]{woo2022etsformer}
G.~Woo, C.~Liu, D.~Sahoo, A.~Kumar, and S.~Hoi.
ETSformer: Exponential smoothing transformers for time-series forecasting.
\emph{arXiv:2202.01381}, 2022.

\bibitem[Zhang and Yan(2023)]{zhang2023crossformer}
Y.~Zhang and J.~Yan.
Crossformer: Transformer utilizing cross-dimension dependency for multivariate
time series forecasting.
In \emph{Proc.\ ICLR}, 2023.

\end{thebibliography}

\end{document}